\documentclass[aps,preprint]{revtex4}%
\usepackage{amsfonts}
\usepackage{amsmath}
\usepackage{amssymb}
\usepackage{graphicx}%
\setcounter{MaxMatrixCols}{30}
%TCIDATA{OutputFilter=latex2.dll}
%TCIDATA{Version=4.10.0.2363}
%TCIDATA{CSTFile=revtex4.cst}
%TCIDATA{Created=Monday, May 24, 2004 10:16:13}
%TCIDATA{LastRevised=Monday, May 24, 2004 10:28:16}
%TCIDATA{<META NAME="GraphicsSave" CONTENT="32">}
%TCIDATA{<META NAME="DocumentShell" CONTENT="Articles\SW\REVTeX 4">}
\newtheorem{theorem}{Theorem}
\newtheorem{acknowledgement}[theorem]{Acknowledgement}
\newtheorem{algorithm}[theorem]{Algorithm}
\newtheorem{axiom}[theorem]{Axiom}
\newtheorem{claim}[theorem]{Claim}
\newtheorem{conclusion}[theorem]{Conclusion}
\newtheorem{condition}[theorem]{Condition}
\newtheorem{conjecture}[theorem]{Conjecture}
\newtheorem{corollary}[theorem]{Corollary}
\newtheorem{criterion}[theorem]{Criterion}
\newtheorem{definition}[theorem]{Definition}
\newtheorem{example}[theorem]{Example}
\newtheorem{exercise}[theorem]{Exercise}
\newtheorem{lemma}[theorem]{Lemma}
\newtheorem{notation}[theorem]{Notation}
\newtheorem{problem}[theorem]{Problem}
\newtheorem{proposition}[theorem]{Proposition}
\newtheorem{remark}[theorem]{Remark}
\newtheorem{solution}[theorem]{Solution}
\newtheorem{summary}[theorem]{Summary}
\newenvironment{proof}[1][Proof]{\noindent\textbf{#1.} }{\ \rule{0.5em}{0.5em}}
\begin{document}
\preprint{HEP/123-qed}
\title[Short title for running header]{Long title that appears on the title page}
\author{First Author}
\affiliation{My Institution}
\author{Second Author}
\affiliation{My Institution}
\author{Third Author}
\affiliation{Other Institution}
\keywords{one two three}
\pacs{PACS number}

\begin{abstract}
Shell document for REV\TeX{} 4.

\end{abstract}
\volumeyear{year}
\volumenumber{number}
\issuenumber{number}
\eid{identifier}
\date[Date text]{date}
\received[Received text]{date}

\revised[Revised text]{date}

\accepted[Accepted text]{date}

\published[Published text]{date}

\startpage{101}
\endpage{102}
\maketitle
\tableofcontents


\section{ Introduction \label{intro}}

Independent Particle States (IPS) have proved a mainstay of molecular
multi-particle quantum mechanics both in physical modelling and numerical
applications for many decades. The major attraction of these states lies in
the conceptually concrete and chemically understandable pictures of molecular
structure and properties that they provide. Apart from an initial choice of
$s$ atomic orbitals and possibly a decision about the molecular point group
and spin symmetry one can obtain an unbiased description of the electronic
states of a molecule by finding an IPS\ that minimizes the energy via the Self
Consistent Field (SCF)\ equations. These equations provide a set of one
electron states $\left\{  \left.  \left\vert \varphi_{j}\right\rangle
\right\vert 1\leq j\leq2s\right\}  $, that are eigenstates of the self
consistent Fock operator corresponding to the eigenvalues $\left\{  \left.
\varepsilon_{j}\right\vert 1\leq j\leq2s\right\}  $ and are unique up to
unitary transformations that mix degenerate eigenstates. If the IPS\ describes
$M$ electrons then in general it is determined by the one-electron eigenstates
corresponding to the $M$ lowest eigenvalues $\left\{  \left.  \varepsilon
_{j}\right\vert 1\leq j\leq M\right\}  $. The wave functions of each of these
one electron states can then be analyzed in terms of the basis of atomic
orbitals, producing the familiar picture of one electronic probability
distributions labelled by one electron energies, that has been invaluable in
the understanding of chemical processes. The physical meaning of the energies
$\left\{  \left.  \varepsilon_{j}\right\vert 1\leq j\leq M\right\}  $ can be
deduced by considering the electron propagator \cite{ortiz} which shows that
they correspond to ionization potentials.

It is often the case, however, that the IPS\ model is not sufficiently
accurate to predict or explain the finer details of chemical reaction
mechanisms or spectroscopic observations. In order to achieve this one must
take into account electron-electron correlation. Traditionally this is
accomplished by perturbational or configurational interaction methods that
increase numerical accuracy at the cost of a clear, tangible, chemical model
of the molecular system. By necessity such models need be based on one
particle properties, one of the motivations of Density Functional Theory, but
in general one particle properties do not determine correlated $N$-electron states.

We have shown \cite{Weiner & Ortiz 2003} that the AGP state for even number
and the GAGP state for an odd number of particles generalize the IPS model and
explicitly contain a description of correlation effects, while are the
\emph{most general} states that produces a \emph{one particle} theory of
electrons. In the IPS\ case the $2N$-electron state is completely determined
by the occupation numbers of the First Order Reduced Density Operator (FORDO)
and a set of \emph{one particle} Natural General Spin Orbitals (NGSOs). The
occupation numbers in this case are equal to one or zero. In the AGP\ case the
$2N$-electron state is completely determined by a set of \emph{one particle}
Canonical General Spin Orbitals (CGSOs) and occupation numbers of the FORDO,
but now the occupation numbers can be any value in the range $\left[
0,1\right]  ,$ though they must be at least doubly degenerate. This theory
explicitly \emph{contains correlation} and generalizes the one-particle theory
based on uncorrelated Independent Particle States (IPS), which it contains as
a special case. The AGP\ model satisfies the requirement that it does produce
a simple pictorial interpretation in terms of occupation numbers and orbitals,
which totally determine the $N$-electron state, and thus can be used as
effective and rigorous chemical modelling tools.

It is important to note the distinction between

\begin{itemize}
\item AGP-NGSO's, which together with a set of doubly degenerate occupation
numbers, determine an AGP-FORDO, that correspond to a \emph{set }of
AGP\ states and

\item AGP-CGSO's, which together with a set of doubly degenerate occupation
numbers, determine a \emph{unique} AGP state.
\end{itemize}

The CGSO's are the NGSO's of the FORDO of a given AGP state, but the NGSO's of
the FORDO\ of this AGP state are \emph{not }in general the CSGO's of this AGP\ state.

In this article we describe how the manifold of $2N$-electron AGP states can
be partitioned into equivalence classes, $\left[  D^{1}\right]  _{2N}$,
indexed by FORDO's $D^{1}$, that are at least doubly degenerate, and how the
elements of each class can be parameterized by a vector
\begin{equation}
\mathbf{\theta\,}=\left(  0,\theta_{2},\cdots,\theta_{s}\right)
,\quad0<\theta_{j}\leq2\pi;\quad2\leq j\leq s,
\end{equation}
and a representative FORDO $D_{0}^{1}.$ Each representative $D_{0}^{1}$ can
parameterized by a vector of geminal occupation numbers%
\begin{equation}
\mathbf{n}=\left(  n_{1},\cdots,n_{s}\right)  ,\quad0\leq n_{j}\leq
1;\quad1\leq j\leq s;\quad%
%TCIMACRO{\tsum \limits_{1\leq j\leq s}}%
%BeginExpansion
{\textstyle\sum\limits_{1\leq j\leq s}}
%EndExpansion
n_{j}=1
\end{equation}
and a representative set of NGSO's $\left\vert \mathbf{\varphi}\right\rangle
=\left\{  \left.  \left\vert \varphi_{j}\right\rangle \right\vert 1\leq
j\leq2s\right\}  $. The energy functional $\mathcal{E}\left(  g^{N}\right)  $
over the manifold of AGP\ states can thus be formulated as a functional
$\mathcal{E}\left(  \mathbf{n},\mathbf{\theta,}\left\vert \mathbf{\varphi
}\right\rangle \right)  $ and we show that the necessary conditions for the
minimization of the energy wrt $\left\vert \mathbf{\varphi}\right\rangle $,
subject to the \emph{normality} constraint, for fixed $\left(  \mathbf{n}%
,\mathbf{\theta}\right)  $ can be formulated in terms of a generalized Fock
operator $F\left(  \mathbf{n},\mathbf{\theta}\right)  $ (we show that is
sufficient to hold the normality constraint in order to obtain
orthonormality). By expressing the total energy in terms of the eigenvalues of
$F\left(  \mathbf{n},\mathbf{\theta}\right)  $ we deduce that the same aufbau
principle as in the IPS\ case should hold i.e. in general the GSO's that are
eigenvectors of $F\left(  \mathbf{n},\mathbf{\theta}\right)  $ with with the
lowest eigenvalues should be associated with the largest occupation numbers.
If the occupation numbers are more than doubly degenerate or if the
eigenvalues of $F\left(  \mathbf{n},\mathbf{\theta}\right)  $ are degenerate
than multiple assignments may have to be examined. Unlike the IPS\ case the
phases of the eigenvectors of $F\left(  \mathbf{n},\mathbf{\theta}\right)  $
\emph{are }important as they effect the optimal value of $\mathbf{\theta}$.
These phases are not determined by the diagonalization of $F\left(
\mathbf{n},\mathbf{\theta}\right)  $ and thus must be determined by an
independent variation. The variational procedure can thus be schematically
displayed as The GSO's and $\mathbf{\theta}$ that are self consistent
solutions to the above scheme determine the CGSO's and together with the self
consistent occupation numbers $\mathbf{n}$ determine the optimal geminal and
$2N$-electron AGP.

The structure of this article is as follows: in section \ref{AGPstate} we
describe the canonical form of a geminal, and thus the AGP\ state, in terms of
geminal occupation numbers $\left\{  \left.  n_{j}\right\vert 1\leq j\leq
s\right\}  $ and CGSOs $\left\{  \left.  \left\vert \psi_{j}\right\rangle
\right\vert 1\leq j\leq2s\right\}  .$ Together with the AGP-FORDO occupation
numbers $\left\{  \left.  N_{j}\right\vert 1\leq j\leq s\right\}  ,$ which are
in $1-1$ correspondence with the geminal occupation numbers, these CGSOs
determine both the AGP\ state and the AGP-FORDO. In section \ref{transform} we
examine the $1-1$ transformations between the AGP $\left\{  \left.
N_{j}\right\vert 1\leq j\leq s\right\}  $ and geminal $\left\{  \left.
n_{j}\right\vert 1\leq j\leq s\right\}  $ occupation numbers. In section
\ref{energy} we discuss the energy expressions in terms of the occupation
numbers $\left\{  \left.  n_{j}\right\vert 1\leq j\leq s\right\}  $, CGSOs
$\left\{  \left.  \left\vert \psi_{j}\right\rangle \right\vert 1\leq
j\leq2s\right\}  $ and equivalence classes $\left[  D^{1}\right]  _{2N}.$ In
section \ref{Evariation} we investigate the necessary conditions needed to
minimize the energy over a manifold of AGP\ states, and derive a generalized
Fock operator that produces the derivative of the energy functional with
respect to NGSO's for fixed FORDO occupation numbers. In section
\ref{procedure} we describe the energy variation procedure in terms of an
iterative process involving the generalized Fock operator, variation over the
equivalence classes $\left[  D^{1}\right]  _{2N}$ and the occupation numbers
$\left\{  \left.  n_{j}\right\vert 1\leq j\leq s\right\}  .$ We conclude in
section \ref{Discussion} with a review of the material presented and a
discussion of its implications.

\section{The $2N$-electron AGP state\label{AGPstate}}

In this section we review the canonical form of a geminal and AGP state and
their associated FORDO's. We also characterize equivalence classes of AGP
states that have the same FORDO.

\subsection{Canonical Form}

The $2N$-electron AGP state, $\left\vert g^{N}\right\rangle $, is the $N^{th}$
antisymmetric power of the $2$-electron state described by the geminal
\begin{equation}
\left\vert g\right\rangle =\sum_{1\leq j\leq s}n_{j}^{\frac{1}{2}}\left\vert
\psi_{j}\psi_{j+s}\right\rangle , \label{geminal}%
\end{equation}
where $\left\{  \left.  n_{j}^{\frac{1}{2}}\right\vert 1\leq j\leq s\right\}
$ are the positive square roots of $\left\{  \left.  n_{j}\right\vert 1\leq
j\leq s\right\}  ,$ and is given by \cite{weiner&ortiz2004}
\begin{equation}
\left\vert g^{N}\right\rangle =\left(  S_{N}\right)  ^{-\frac{1}{2}}%
\sum_{1\leq j_{1}<\cdots<j_{N}\leq s}n_{j_{1}}^{\frac{1}{2}}\cdots n_{j_{s}%
}^{\frac{1}{2}}\left\vert \psi_{j_{1}}\psi_{j_{1}+s}\cdots\psi_{j_{N}}%
\psi_{j_{N}+s}\right\rangle . \label{agp}%
\end{equation}
The normalization factor, $\left(  S_{N}\right)  ^{-\frac{1}{2}}$, is given in
terms of the elementary symmetric function
\begin{equation}
S_{N}=\sum_{1\leq j_{1}<\cdots<j_{N}\leq s}n_{j_{1}}\cdots n_{j_{N}},
\end{equation}
and the geminals $\left\vert g\right\rangle $ represent normalized two
electron states so that
\begin{equation}
\sum_{1\leq j\leq s}n_{j}=1.
\end{equation}
It is evident from the expansion eq. $\left(  \ref{geminal}\right)  $ that the
geminal and geminal power state are at the least invariant to the group
\begin{equation}
SU\left(  2\right)  ^{s}=\overset{s\text{-factors}}{\overbrace{\,SU\left(
2\right)  \times\cdots\times SU\left(  2\right)  }}%
\end{equation}
that is formed by $2\times2$ \emph{special} unitary transformations of the
paired CGSO's $\left\{  \left.  \left\vert \psi_{j}\right\rangle ,\left\vert
\psi_{j+s}\right\rangle \right\vert 1\leq j\leq s\right\}  ,$ while if
degeneracies exist amongst the $\left\{  \left.  n_{j}\right\vert 1\leq j\leq
s\right\}  $ this invariance group is larger \cite{Weiner2004}. Hence the
canonical CSGO's are unique up to transformations belonging this invariance group.

The CGSO's $\left\{  \left.  \left\vert \psi_{j}\right\rangle \right\vert
1\leq j\leq2s\right\}  $ can be expanded in terms of a spin orbitals basis
$\left\{  \left.  \left\vert \varphi_{j\alpha}\right\rangle ,\left\vert
\varphi_{j\beta}\right\rangle \right\vert 1\leq j\leq s\right\}  $ of one
electron state space $\mathcal{H}^{1}$ as
\begin{equation}
\left\vert \psi_{j}\right\rangle =\sum_{1\leq j\leq s}\left\{  a_{j}\left\vert
\varphi_{j\alpha}\right\rangle +b_{j}\left\vert \varphi_{j\beta}\right\rangle
\right\}  ,
\end{equation}
leading to one electron GSO wavefunctions $\left(  \text{see Appendix
\ref{basis}}\right)  $
\begin{equation}
\psi_{j}\left(  \mathbf{r,}\zeta\right)  =\left\langle \left.  \mathbf{r,}%
\zeta\right\vert \psi_{j}\right\rangle .
\end{equation}
It should be noted that if expansion eq. $\left(  \ref{geminal}\right)  $ is
used to define the canonical expansion it is usually \emph{not possible} to
find a set of CSGO's such that
\begin{equation}
\psi_{j+s}\left(  \mathbf{r,}\zeta\right)  =\psi_{j}\left(  \mathbf{r,}%
\zeta\right)  ^{\ast}.
\end{equation}


\subsection{First Order Reduced Density Operators}

The states $\left\vert g\right\rangle $ and $\left\vert g^{N}\right\rangle $
define FORDO's given by \cite{Coleman}
\begin{equation}
D^{1}\left(  g\right)  =\sum_{1\leq j\leq s}n_{j}\left\{  \left\vert \psi
_{j}\right\rangle \left\langle \psi_{j}\right\vert +\left\vert \psi
_{j+s}\right\rangle \left\langle \psi_{j+s}\right\vert \right\}  \label{d1g}%
\end{equation}
and
\begin{equation}
D_{2N}^{1}=\sum_{1\leq j\leq s}N_{j}\left\{  \left\vert \psi_{j}\right\rangle
\left\langle \psi_{j}\right\vert +\left\vert \psi_{j+s}\right\rangle
\left\langle \psi_{j+s}\right\vert \right\}  \label{dng}%
\end{equation}
respectively. It is clear from the expansions eqs. $\left(  \ref{d1g}\right)
$ and $\left(  \ref{dng}\right)  $ that the CGSO's $\left\{  \left.
\left\vert \psi_{j}\right\rangle \right\vert 1\leq j\leq r\right\}  $ are
simultaneously the Natural General Spin Orbitals (NGSO's) of $D^{1}\left(
g\right)  $ and $D^{1}\left(  g^{N}\right)  $. However it is very important to
note that the NGSO's of $D^{1}\left(  g\right)  $ and $D^{1}\left(
g^{N}\right)  $ \emph{are not }in general the CGSO's of $\left\vert
g\right\rangle $. This is a ramification of

\begin{enumerate}
\item the general invariance of FORDO's to the multiplication of their NGSO's
by complex phase factors
\begin{equation}
D_{2}^{1}=\sum_{1\leq j\leq r}\nu_{j}\left\{  \left\vert \phi_{j}\right\rangle
\left\langle \phi_{j}\right\vert \right\}  =\sum_{1\leq j\leq r}\nu
_{j}\left\{  e^{i\theta_{j}}\left\vert \phi_{j}\right\rangle \left\langle
\phi_{j}\right\vert e^{-i\theta_{j}}\right\}  ,
\end{equation}
where $\left\{  \left.  \nu_{j}\right\vert 1\leq j\leq s\right\}  $ are
occupation numbers and

\item the observation that geminals \emph{are not }invariant to such phase
changes of their CGSO's.
\end{enumerate}

One can form equivalence classes $\left[  D_{2}^{1}\right]  $ of geminals
$\left\{  \left\vert g\left(  \mathbf{\theta}\right)  \right\rangle \right\}
$ that have the same FORDO $D_{2}^{1}$ and equivalence classes $\left[
D_{2N}^{1}\right]  $ of $2N$-electron AGP states $\left\{  \left\vert
g^{N}\left(  \mathbf{\theta}\right)  \right\rangle \right\}  ,$that have the
same FORDO $D_{2N}^{1}$. For a fixed $N$ these classes correspond to each
other in a bijective manner \cite{coleman} and can be parameterized by
$\mathbf{\theta}$, where
\begin{align}
&  \left\vert g\left(  \mathbf{\theta}\right)  \right\rangle =\sum_{1\leq
j\leq s}n_{j}^{\frac{1}{2}}e^{i\theta_{j}}\left\vert \psi_{j}\psi
_{j+s}\right\rangle =\sum_{1\leq j\leq s}n_{j}^{\frac{1}{2}}\left\vert
e^{i\varphi_{j}}\psi_{j}e^{i\varphi_{j+s}}\psi_{j+s}\right\rangle =\sum_{1\leq
j\leq s}n_{j}^{\frac{1}{2}}\left\vert \chi_{j}\chi_{j+s}\right\rangle
\nonumber\\
&  \qquad\qquad\qquad\qquad\qquad\qquad\qquad\qquad\qquad\qquad\qquad
\theta_{1}=0;\,0<\theta_{j}\leq2\pi,2\leq j\leq s,
\end{align}
$\chi_{j}=e^{i\varphi_{j}}\psi_{j},\mathbf{\theta=}\left(  0,\theta_{2}%
,\cdots,\theta_{s}\right)  $ and $\varphi_{j}+\varphi_{j+s}=\theta_{j}$.
Choosing $\theta_{1}=0$ eliminates the overall phase factor indeterminacy of
$\left\vert g\left(  \mathbf{\theta}\right)  \right\rangle $. The individual
phase factors $\left\{  \left.  \varphi_{j}\right\vert 1\leq j\leq s\right\}
$ are not unique, though their sums $\left\{  \left.  \theta_{j}\right\vert
1\leq j\leq s\right\}  $ are, allowing one to obtain a unique representative
set $\left\{  \left.  \varphi_{j}\right\vert 1\leq j\leq s\right\}  $ by
defining
\begin{equation}
\varphi_{j}=\frac{\theta_{j}}{2},\varphi_{j+s}=\frac{\theta_{j}}{2};\quad1\leq
j\leq s.
\end{equation}
Thus one has the correspondence
\begin{equation}
\left\{  \left\vert g\left(  \mathbf{\theta}\right)  \right\rangle
\longleftrightarrow\left\vert g^{N}\left(  \mathbf{\theta}\right)
\right\rangle ;\quad\theta_{1}=0;\,0<\theta_{j}\leq2\pi,2\leq j\leq s\right\}
\rightarrow D_{2N}^{1}.
\end{equation}
The AGP\ states $\left\vert g^{N}\left(  \mathbf{\theta}\right)  \right\rangle
$ are clearly distinct as they have different Second Order Reduced Density
Operators (SORDO's), $D_{2N}^{2}\left(  \mathbf{\theta}\right)  $, given by
\cite{Coleman}
\begin{align}
D_{2N}^{2}\left(  \mathbf{\theta}\right)   &  =\left(  S_{N}\right)
^{-\frac{1}{2}}\left\{  \sum_{1\leq j\leq s}N_{j}\left\vert \psi_{j}\psi
_{j+s}\right\rangle \left\langle \psi_{j}\psi_{j+s}\right\vert +\sum_{1\leq
j,k\leq s}n_{j}^{\frac{1}{2}}n_{k}^{\frac{1}{2}}e^{i\left(  \theta_{j}%
-\theta_{k}\right)  }S_{N-1}\left(  \jmath k\right)  \left\vert \psi_{j}%
\psi_{j+s}\right\rangle \left\langle \psi_{k}\psi_{k+s}\right\vert \right.
\nonumber\\
&  \qquad\qquad\qquad\qquad\qquad\qquad\qquad\qquad\qquad\left.  +\sum_{1\leq
j<k\leq s}n_{j}n_{k}S_{N-2}\left(  \jmath k\right)  \left\vert \psi_{j}%
\psi_{k}\right\rangle \left\langle \psi_{j}\psi_{k}\right\vert \right\}  ,
\label{D2}%
\end{align}
where
\begin{equation}
S_{N-1}\left(  \jmath k\right)  =\sum_{\substack{1\leq j_{1}<\cdots
<j_{N-1}\leq s\\j,k\notin\left\{  j_{1},\cdots,j_{N-1}\right\}  }}n_{j_{1}%
}\cdots n_{j_{N-1}}%
\end{equation}
and
\begin{equation}
S_{N-2}\left(  \jmath k\right)  =\sum_{\substack{1\leq j_{1}<\cdots
<j_{N-2}\leq s\\j,k\notin\left\{  j_{1},\cdots,j_{N-2}\right\}  }}n_{j_{1}%
}\cdots n_{j_{N-2}}.
\end{equation}
Note that $S_{N-1}\left(  \jmath\jmath\right)  =S_{N-2}\left(  \jmath
\jmath\right)  =0.$ The SORDO eq. $\left(  \ref{D2}\right)  $ can be expanded
in terms of $\mathbf{\theta}$ dependent CGSO's $\left\{  \left.  \left\vert
e^{i\varphi_{j}}\psi_{j}\right\rangle ,\left\vert e^{i\varphi_{j+s}}\psi
_{j+s}\right\rangle \right\vert 1\leq j\leq s\right\}  ,$ where $\left\{
\left.  \varphi_{j}+\varphi_{j+s}=\theta_{j}\right\vert 1\leq j\leq s\right\}
$ as
\begin{align}
&  D_{2N}^{2}\left(  \mathbf{\theta}\right)  =\left(  S_{N}\right)
^{-\frac{1}{2}}\left\{  \sum_{1\leq j\leq s}N_{j}\left\vert e^{i\varphi_{j}%
}\psi_{j}e^{i\varphi_{j+s}}\psi_{j+s}\right\rangle \left\langle e^{i\varphi
_{j}}\psi_{j}e^{i\varphi_{j+s}}\psi_{j+s}\right\vert \right. \nonumber\\
&  \qquad\qquad\qquad\qquad+\sum_{1\leq j\neq k\leq s}n_{j}^{\frac{1}{2}}%
n_{k}^{\frac{1}{2}}S_{N-1}\left(  \jmath k\right)  \left\vert e^{i\varphi_{j}%
}\psi_{j}e^{i\varphi_{j+s}}\psi_{j+s}\right\rangle \left\langle e^{i\varphi
_{k}}\psi_{k}e^{i\varphi_{k+s}}\psi_{k+s}\right\vert \nonumber\\
&  \qquad\qquad\qquad\qquad\qquad\qquad\qquad\qquad\left.  +\sum_{1\leq
j<k\leq s}n_{j}n_{k}S_{N-2}\left(  \jmath k\right)  \left\vert e^{i\varphi
_{j}}\psi_{j}e^{i\varphi_{j+s}}\psi_{k}\right\rangle \left\langle
e^{i\varphi_{j}}\psi_{j}e^{i\varphi_{j+s}}\psi_{k}\right\vert \right\}  ,
\label{en2}%
\end{align}
giving
\begin{align}
&  D^{2}\left(  g^{N}\left(  \mathbf{\theta}\right)  \right)  =\left(
S_{N}\right)  ^{-\frac{1}{2}}\left\{  \sum_{1\leq j\leq s}N_{j}\left\vert
\chi_{j}\chi_{j+s}\right\rangle \left\langle \chi_{j}\chi_{j+s}\right\vert
\right. \nonumber\\
&  \qquad\qquad\qquad\qquad+\sum_{1\leq j\neq k\leq s}n_{j}^{\frac{1}{2}}%
n_{k}^{\frac{1}{2}}S_{N-1}\left(  \jmath k\right)  \left\vert \chi_{j}%
\chi_{j+s}\right\rangle \left\langle \chi_{k}\chi_{k+s}\right\vert \nonumber\\
&  \qquad\qquad\qquad\qquad\qquad\qquad\qquad\qquad\left.  +\sum_{1\leq
j<k\leq s}n_{j}n_{k}S_{N-2}\left(  \jmath k\right)  \left\vert \chi_{j}%
\chi_{k}\right\rangle \left\langle \chi_{j}\chi_{k}\right\vert \right\}  ,
\end{align}
where $\left\vert \chi_{j}\right\rangle =\left\vert e^{i\varphi_{j}}\psi
_{j}\right\rangle .$

\section{The correspondence between AGP and Geminal Occupation
Numbers\label{transform}}

Any set of $s$ positive numbers $\left\{  \left.  N_{j}\right\vert 1\leq j\leq
s\right\}  $ that have values between $0$ and $1$ and add up to $N,$ can be
interpreted as the occupation numbers of a FORDO that is at least doubly
degenerate and thus can be associated with an AGP\ state. These numbers are in
$1-1$ correspondence with another set of positive numbers $\left\{  \left.
n_{j}\right\vert 1\leq j\leq s\right\}  ,$ that have values between $0$ and
$1$ and add up to $1$, which can be interpreted as the occupation numbers of a
FORDO\ of a geminal. The geminal occupation numbers $\left\{  \left.
n_{j}\right\vert 1\leq j\leq s\right\}  $ together with an arbitrary set of
$2s$ orthonormal one electron Canonical General Spin Orbitals (CGSO's) then
\emph{uniquely }determine a $2N$ electron AGP state. Any state that has a
FORDO that has at least doubly degenerate occupation numbers is covered by a
set of AGP\ states i.e. there exists AGP\ states that have the same FORDO. In
most applications e.g. expectation values of observables, it is sufficient to
know the $2N$ electron state occupation numbers in terms of the geminal
occupation numbers, which serve as a parameterization of the state occupation
numbers, and we review this relationship first. The inverse function i.e.
geminal occupation numbers in terms of state occupation numbers are
significant in designing possible AGP\ states and we present a numerical
construction of this function.

\subsection{Transformation of Geminal to AGP Occupation Numbers}

The vector of occupation numbers $\mathbf{N}=\left(  N_{1},\cdots
,N_{s}\right)  $ of an AGP\ state $\left\vert g^{N}\right\rangle $ can be
expressed as a bijective function, $\mathcal{N}\left(  \mathbf{n}\right)  $,
of the geminal occupation numbers $\mathbf{n}=\left(  n_{1},\cdots
,n_{s}\right)  ,$ where
\begin{equation}
\mathcal{N}:\mathfrak{D}_{N}\subset\mathbb{R}^{s}\rightarrow\mathfrak{D}%
_{n}\subset\mathbb{R}^{s}.
\end{equation}
The domain $\mathfrak{D}_{N}$ and range $\mathfrak{D}_{n}$ are given by%
\begin{align}
\mathfrak{D}_{N}  &  =\left(  N_{1},\cdots,N_{s}\right)  ;\quad0\leq N_{j}%
\leq1,\quad1\leq j\leq s\nonumber\\
\sum_{1\leq j\leq s}N_{j}  &  =N
\end{align}
and
\begin{align}
\mathfrak{D}_{n}  &  =\left(  n_{1},\cdots,n_{s}\right)  ;\quad0\leq n_{j}%
\leq1,\quad1\leq j\leq s\nonumber\\
\sum_{1\leq j\leq s}n_{j}  &  =1,
\end{align}
where
\begin{equation}
N_{j}=\mathcal{N}_{j}\left(  \mathbf{n}\right)  =n_{j}\frac{S_{N-1}\left(
\jmath\right)  }{S_{N}};\quad1\leq j\leq s
\end{equation}
and
\begin{equation}
S_{N-1}\left(  \jmath\right)  =\sum_{\substack{1\leq j_{1}<\cdots<j_{N-1}\leq
s\\j\notin\left\{  j_{1},\cdots,j_{N-1}\right\}  }}n_{j_{1}}\cdots n_{j_{N-1}%
}.
\end{equation}


\subsection{Transformation of AGP to Geminal Occupation Numbers}

The inverse function $\mathcal{N}^{-1}\left(  \mathbf{N}\right)  $ is not
explicitly known, but one can obtain a linear approximation that suffices in
practical applications to numerical determine values of $\mathbf{n}$ from
known values of $\mathbf{N}$ i.e. the determination of $\mathbf{n}_{0}%
+\delta\mathbf{n}$ from $\mathbf{N}_{0}+\delta\mathbf{N}$, where
$\mathbf{N}_{0}$ is a known initial vector of occupation numbers of the
$2N$-electron AGP state e.g. an IPS,\ where
\begin{equation}
\mathbf{N}_{0}=\left(  \overset{N}{\overbrace{1,\cdots,1}},\overset
{s-N}{\overbrace{0,\cdots,0}}\right)
\end{equation}
and
\begin{equation}
\mathbf{n}_{0}=\left(  \overset{N}{\overbrace{\frac{1}{N},\cdots,\frac{1}{N}}%
},\overset{s-N}{\overbrace{0,\cdots,0}}\right)  .
\end{equation}
The Taylor series expansion
\begin{equation}
\mathcal{N}^{-1}\left(  \mathbf{N}_{0}+\delta\mathbf{N}\right)  =\mathcal{N}%
^{-1}\left(  \mathbf{N}_{0}\right)  +d^{1}\mathcal{N}^{-1}\left(
\mathbf{N}_{0}\right)  \left(  \delta\mathbf{n}\right)  +\cdots,
\end{equation}
where
\begin{equation}
d^{1}\mathcal{N}^{-1}\left(  \mathbf{N}_{0}\right)  =\left[  d^{1}%
\mathcal{N}\left(  \mathbf{n}_{0}\right)  \right]  ^{-1}%
\end{equation}
and the Jacobian of the transformation $\mathcal{N}$ at the point
$\mathbf{n}_{0}$ is
\begin{align}
\left[  d^{1}\mathcal{N}\left(  \mathbf{n}_{0}\right)  \right]  _{jk}  &
=\frac{\partial\mathcal{N}_{j}\left(  \mathbf{n}_{0}\right)  }{\partial n_{k}%
}=\frac{\partial}{\partial n_{k}}\left\{  n_{j}\frac{S_{N-1}\left(
\jmath\right)  }{S_{N}}\right\}  =\frac{1}{S_{N}}\left(  \delta_{jk}%
S_{N-1}\left(  \jmath\right)  +n_{j}\frac{\partial S_{N-1}\left(
\jmath\right)  }{\partial n_{k}}-n_{j}\frac{S_{N-1}\left(  \jmath\right)
}{S_{N}}\frac{\partial S_{N}}{\partial n_{k}}\right) \nonumber\\
&  =\frac{1}{S_{N}}\left(  \delta_{jk}S_{N-1}\left(  \jmath\right)
+n_{j}S_{N-2}\left(  \jmath k\right)  -n_{j}\frac{1}{S_{N}}S_{N-1}\left(
\jmath\right)  S_{N-1}\left(  k\right)  \right)
\end{align}
Hence the Jacobian of the inverse transformation $D^{1}\mathcal{N}^{-1}\left(
\mathbf{N}_{0}\right)  $ at the point $\mathbf{N}_{0}$ is given by the inverse
of the matrix $\left\{  \left.  \left[  d^{1}\mathcal{N}\left(  \mathbf{n}%
_{0}\right)  \right]  _{jk}\right\vert 1\leq j,k\leq s\right\}  $ i.e.
\begin{equation}
d^{1}\mathcal{N}^{-1}\left(  \mathbf{N}_{0}\right)  =\left[  d^{1}%
\mathcal{N}\left(  \mathbf{n}_{0}\right)  \right]  ^{-1}%
\end{equation}
and the incremental change, $\delta\mathbf{n}$ of $\mathbf{n}_{0}$ is
\begin{equation}
d^{1}\mathcal{N}^{-1}\left(  \mathbf{N}_{0}\right)  \left(  \delta
\mathbf{N}\right)  _{j}=\sum_{1\leq k\leq s}\left[  d^{1}\mathcal{N}\left(
\mathbf{n}_{0}\right)  \right]  _{jk}^{-1}\left(  \delta\mathbf{N}\right)
_{k}.
\end{equation}
We can then construct a piecewise linear fit to $\mathcal{N}^{-1}$ as%
\begin{equation}
\mathbf{n}\simeq\mathbf{\tilde{n}}_{m}=\mathcal{N}^{-1}\left(  \mathbf{N}%
\right)  =\mathbf{n}_{0}+\sum_{1\leq k\leq m}\left[  d^{1}\mathcal{N}\left(
\mathbf{n}_{0}+\sum_{1\leq j\leq k-1}\left[  d^{1}\mathcal{N}\left(
\mathbf{n}_{0}\right)  \right]  ^{-1}\left(  \delta\mathbf{N}_{j}\right)
\right)  \right]  ^{-1}\left(  \delta\mathbf{N}_{k}\right)  ,
\end{equation}
where
\begin{equation}
\mathbf{N=N}_{0}+\sum_{1\leq j\leq k-1}\delta\mathbf{N}_{j}%
\end{equation}
and
\begin{equation}
\sum_{1\leq k\leq s}\delta\mathbf{N}_{jk}=0.
\end{equation}
See Appendix \label{linfit}. We can check the accuracy of the fit by
evaluating
\begin{equation}
\mathbf{\Delta}_{m}=\mathcal{N}\left(  \mathbf{\tilde{n}}_{m}\right)
-\mathbf{N}_{m}%
\end{equation}
and checking if
\begin{equation}
\left\vert \mathbf{\Delta}_{mj}\right\vert \leq\varepsilon_{j};\quad1\leq
j\leq s
\end{equation}
If the error is too big one can then decrease the size of $\delta
\mathbf{N}_{m}$.

\section{Energy Functionals \label{Energy}}

In this section we first describe the energy of a $2N$-electron AGP state in
terms of CGSO's and geminal occupation numbers then in terms of a density
matrix functional $\mathcal{F}\left(  D^{1}\right)  $ defined on the
equivalence classes $\left[  D_{2N}^{1}\right]  $. The functional
$\mathcal{F}\left(  D^{1}\right)  $ is significant in two aspects, it is
essential in our proposed optimization scheme in terms of generalized Fock
operators and it can serve as a rigorous explicitly known model in Density
Matrix Functional Theory (DMFT) \cite{DMFT}.

\subsection{Energy in terms of Geminal Occupation Numbers and CGSO's}

The energy, $\mathcal{E}\left(  \mathbf{n,\psi}\right)  ,$ of a 2$N$-electron
AGP\ state for a molecule at a fixed nuclear configuration $\left(
\mathbf{R}_{1},\cdots,\mathbf{R}_{N_{nuc}}\right)  $ in terms of the CGSO's
amplitudes $\left\{  \left.  \left\langle \left.  \mathbf{r,\xi}\right\vert
\psi_{j}\right\rangle \right\vert 1\leq j\leq r\right\}  $ and geminal
occupation numbers $\mathbf{n}$ is
\begin{align}
\mathcal{E}\left(  \mathbf{n,\psi}\right)  =  &  \frac{\mathcal{H}\left(
\mathbf{n,\psi}\right)  }{S_{N}\left(  \mathbf{n}\right)  }=\frac{1}%
{S_{N}\left(  \mathbf{n}\right)  }\left\{  \sum_{1\leq j\leq s}N_{j}\left\{
\left\langle \psi_{j}\left\vert h_{1}\psi_{j}\right.  \right\rangle
+\left\langle \psi_{j+s}\left\vert h_{1}\psi_{j+s}\right.  \right\rangle
+\left\langle \psi_{j}\psi_{j+s}|h_{2}|\psi_{j}\psi_{j+s}\right\rangle
\right\}  \right. \nonumber\\
&  +\sum_{1\leq j,k\leq s}n_{j}^{\frac{1}{2}}n_{k}^{\frac{1}{2}}S_{N-1}\left(
\jmath k\right)  \left\langle \psi_{j}\psi_{j+s}|h_{2}|\psi_{k}\psi
_{k+s}\right\rangle \left.  +\sum_{1\leq i,j\leq s}n_{i}n_{j}S_{N-2}\left(
\jmath k\right)  \left\langle \psi_{j}\psi_{k}|||\psi_{j}\psi_{k}\right\rangle
\right\}  . \label{energy}%
\end{align}
$\lceil$Note: Mathematically the energy functional should really be denoted as
$\mathcal{E}\left(  \mathbf{n,\psi}^{\ast}\mathbf{,\psi}\right)  $ as
$\mathbf{\psi}^{\ast}$ and $\mathbf{\psi}$ are treated as independent
variables even though $\mathbf{\psi}$ determines $\mathbf{\psi}^{\ast}%
$.$\rfloor$ The one-electron matrix elements $\left\{  \left\langle \psi
_{j}\left\vert h_{1}\psi_{k}\right.  \right\rangle \right\}  $ wrt the CGSO's
are
\begin{equation}
\left\langle \psi_{j}\left\vert h_{1}\psi_{k}\right.  \right\rangle =-\int
\psi_{j}^{\ast}\left(  \mathbf{r,\xi}\right)  \left\{  \sum_{_{1\leq l\leq
N_{nuc}}}\frac{Z_{l}}{\left\Vert \mathbf{R}_{l}-\mathbf{r}\right\Vert }%
+\nabla_{\mathbf{r}}^{2}\right\}  \psi_{k}\left(  \mathbf{r,\xi}\right)
d\mathbf{r}d\xi;1\leq j,k\leq2s, \label{h1}%
\end{equation}
the antisymmetric two-electron matrix elements $\left\{  \left\langle
jk||lm\right\rangle \right\}  $
\begin{align}
\left\langle \psi_{j}\psi_{j+s}|h_{2}|\psi_{j}\psi_{j+s}\right\rangle  &
=\int\left\{  \psi_{j}^{\ast}\left(  \mathbf{r}_{1}\mathbf{,\xi}_{1}\right)
\psi_{k}^{\ast}\left(  \mathbf{r}_{2}\mathbf{,\xi}_{2}\right)  \frac
{1}{\left\Vert \mathbf{r}_{1}-\mathbf{r}_{2}\right\Vert }\psi_{l}\left(
\mathbf{r}_{1}\mathbf{,\xi}_{1}\right)  \psi_{m}\left(  \mathbf{r}%
_{2}\mathbf{,\xi}_{2}\right)  \right. \nonumber\\
&  \qquad-\left.  \psi_{j}^{\ast}\left(  \mathbf{r}_{1}\mathbf{,\xi}%
_{1}\right)  \psi_{k}^{\ast}\left(  \mathbf{r}_{2}\mathbf{,\xi}_{2}\right)
\frac{1}{\left\Vert \mathbf{r}_{1}-\mathbf{r}_{2}\right\Vert }\psi_{l}\left(
\mathbf{r}_{2}\mathbf{,\xi}_{1}\right)  \psi_{m}\left(  \mathbf{r}%
_{1}\mathbf{,\xi}_{2}\right)  \right\}  d\mathbf{r}_{1}d\mathbf{r}_{2}d\xi
_{1}d\xi_{2}\nonumber\\
&  \left.  {}\right.  \qquad\qquad\qquad\qquad\qquad\qquad\qquad\qquad
\qquad1\leq j,k\leq2s;1\leq l,m\leq2s \label{h2}%
\end{align}
and $\left\{  \left\langle jk|||jk\right\rangle \right\}  $
\begin{align}
\left\langle \psi_{j}\psi_{k}|||\psi_{j}\psi_{k}\right\rangle  &
=\left\langle jk||jk\right\rangle +\left\langle j\,k+s||j\,k+s\right\rangle
+\left\langle j+s\,k||j+s\,k\right\rangle +\left\langle
j+s\,k+s||j+s\,k+s\right\rangle \nonumber\\
&  \left.  {}\right.  \qquad\qquad\qquad\qquad\qquad\qquad\qquad\qquad
\qquad\qquad\qquad\qquad\qquad\qquad1\leq j,k\leq s.
\end{align}
The integral kernels of the operators $h_{1}$ and $h_{2}$ are
\begin{align}
&  h_{1}\left(  \mathbf{r}_{1}\mathbf{,\xi}_{1};\mathbf{r}_{1}^{\prime
}\mathbf{,\xi}_{1}^{\prime}\right)  =\left\langle \mathbf{r}_{1}\mathbf{,\xi
}_{1}\left\vert \left(  h_{1}\right)  \mathbf{r}_{1}^{\prime}\mathbf{,\xi}%
_{1}^{\prime}\right.  \right\rangle \nonumber\\
&  \qquad\qquad\qquad\qquad\left.  =\right.  \sum_{_{1\leq l\leq N_{nuc}}%
}\frac{Z_{l}}{\left\Vert \mathbf{R}_{l}-\mathbf{r}_{1}\right\Vert }%
\delta\left(  \mathbf{r}_{1}-\mathbf{r}_{1}^{\prime}\right)  \delta\left(
\mathbf{\xi}_{1}-\mathbf{\xi}_{1}^{\prime}\right)  +\nabla_{\mathbf{r}_{1}%
}^{2}\delta\left(  \mathbf{r}_{1}-\mathbf{r}_{1}^{\prime}\right) \nonumber\\
&  h_{2}\left(  \mathbf{r}_{1}\mathbf{,\xi}_{1},\mathbf{r}_{2}\mathbf{,\xi
}_{2};\mathbf{r}_{1}^{\prime}\mathbf{,\xi}_{1}^{\prime},\mathbf{r}_{2}%
^{\prime}\mathbf{,\xi}_{2}^{\prime}\right)  =\left\langle \mathbf{r}%
_{1}\mathbf{,\xi}_{1},\mathbf{r}_{2}\mathbf{,\xi}_{2}\left.  \left(
h_{2}\right)  \mathbf{r}_{1}^{\prime}\mathbf{,\xi}_{1}^{\prime},\mathbf{r}%
_{2}^{\prime}\mathbf{,\xi}_{2}^{\prime}\right\vert \right\rangle \nonumber\\
&  \qquad\qquad\qquad\qquad\left.  =\right.  \frac{1}{\left\Vert
\mathbf{r}_{1}-\mathbf{r}_{2}\right\Vert }\left\{  \delta\left(
\mathbf{r}_{1}-\mathbf{r}_{1}^{\prime}\right)  \delta\left(  \mathbf{r}%
_{2}-\mathbf{r}_{2}^{\prime}\right)  \delta\left(  \mathbf{\xi}_{1}%
-\mathbf{\xi}_{1}^{\prime}\right)  \delta\left(  \mathbf{\xi}_{2}-\mathbf{\xi
}_{2}^{\prime}\right)  \right. \nonumber\\
&  \qquad\qquad\qquad\qquad\qquad\qquad\qquad\qquad-\left.  \delta\left(
\mathbf{r}_{1}-\mathbf{r}_{2}^{\prime}\right)  \delta\left(  \mathbf{r}%
_{2}-\mathbf{r}_{1}^{\prime}\right)  \delta\left(  \mathbf{\xi}_{1}%
-\mathbf{\xi}_{2}^{\prime}\right)  \delta\left(  \mathbf{\xi}_{2}-\mathbf{\xi
}_{1}^{\prime}\right)  \right\}
\end{align}


\subsection{Energy in terms of CGSO's Phases and FORDO's}

In terms of $\theta$ dependent CGSO's $\left\{  \left\vert e^{i\varphi_{j}%
}\psi_{j}\right\rangle ;1\leq j\leq2s\right\}  $ defined in eq. $\left(
\ref{en2}\right)  $ we obtain an energy expression factored through the
equivalence classes $\left[  D^{1}\right]  _{2N}$
\begin{align}
\mathcal{E}\left(  \mathbf{n},\mathbf{\psi},\mathbf{\theta}\right)  =  &
\frac{\mathcal{H}\left(  \mathbf{n},\mathbf{\psi},\mathbf{\theta}\right)
}{S_{N}\left(  \mathbf{n}\right)  }=\frac{1}{S_{N}\left(  \mathbf{n}\right)
}\left\{  \sum_{1\leq j\leq s}N_{j}\left\{  \left\langle \psi_{j}\left\vert
h_{1}\psi_{j}\right.  \right\rangle +\left\langle \psi_{j+s}\left\vert
h_{1}\psi_{j+s}\right.  \right\rangle +\left\langle \psi_{j}\psi_{j+s}%
|h_{2}|\psi_{j}\psi_{j+s}\right\rangle \right\}  \right. \nonumber\\
&  +\sum_{1\leq j,k\leq s}n_{j}^{\frac{1}{2}}n_{k}^{\frac{1}{2}}e^{i\left(
\theta_{j}-\theta_{k}\right)  }S_{N-1}\left(  \jmath k\right)  \left\langle
\psi_{j}\psi_{j+s}|h_{2}|\psi_{k}\psi_{k+s}\right\rangle \nonumber\\
&  \qquad\qquad\qquad\qquad\qquad\qquad\qquad\qquad\qquad\qquad\left.
+\sum_{1\leq i,j\leq s}n_{i}n_{j}S_{N-2}\left(  \jmath k\right)  \left\langle
\psi_{j}\psi_{k}|||\psi_{j}\psi_{k}\right\rangle \right\}  . \label{thetaen}%
\end{align}
where $\left\vert \mathbf{\psi}\right\rangle $ is a representative set of
CGSO's for the class $\left[  D^{1}\right]  _{2N}$. As $\left(  \mathbf{n}%
,\mathbf{\psi}\right)  $ uniquely defines a $2N$ -electron AGP state
FORDO\ the expression eq. $\left(  \ref{thetaen}\right)  $ can be used to
define a density matrix energy functional
\begin{equation}
\mathcal{F}\left(  D^{1}\right)  =\min_{\mathbf{\theta\in}\Theta}%
\mathcal{E}\left(  \mathbf{n},\mathbf{\psi},\mathbf{\theta}\right)
\end{equation}
over the manifold of AGP\ states. Noting that $D^{1}=D^{1}\left(
\mathbf{n},\mathbf{\psi}\right)  $ this functional can then be minimized to
give
\begin{equation}
\min_{D^{1}\in\mathcal{S}_{1}}\mathcal{F}\left(  D^{1}\right)  =\min
_{\mathbf{n\in}\mathcal{D}_{n};\left\vert \mathbf{\psi}\right\rangle
\in\mathcal{C}}\mathcal{F}\left(  D^{1}\left(  \mathbf{n},\mathbf{\psi
}\right)  \right)  =\min_{\left\vert g\right\rangle \in\mathcal{H}^{2}%
}\mathcal{E}\left(  \left\vert g^{N}\right\rangle \right)
\end{equation}
which characterizes the $2N$-electron AGP state that minimizes the total energy.

\section{Necessary Conditions for the Minimization of the Energy in terms of
Geminal Occupation Numbers, and $\theta$ dependent CGSO's \label{Evariation}}

In order to solve the constrained optimization problem
\begin{equation}
\min_{\left(  \mathbf{n,\psi,\theta}\right)  }\mathcal{E}\left(
\mathbf{n,\psi,\theta}\right)  \label{min}%
\end{equation}
subject to the constraints%
\begin{equation}
\theta_{1}=0;\,0<\theta_{j}\leq2\pi,2\leq j\leq s \label{thetacon}%
\end{equation}%
\begin{align}
\left\langle \left\vert \psi_{j}\right\vert \psi_{k}\right\rangle  &
=\delta_{jk};\quad1\leq j,k\leq2s\label{norm}\\
\mathbf{\psi}  &  \mathbf{\in}\mathcal{R}_{\left[  D_{N}^{1}\right]  }
\label{rep}%
\end{align}
and
\begin{subequations}
\begin{align}
&  0\leq n_{j}\leq1;\quad1\leq j\leq s\label{pos1}\\
&  \sum_{1\leq j\leq s}n_{j}=1
\end{align}
we utilize the method of Lagrange multipliers based on the Lagrangian
\end{subequations}
\begin{equation}
\mathcal{\tilde{E}}\left(  \mathbf{n,\psi,\theta}\right)  =\mathcal{E}\left(
\mathbf{n,\psi,\theta}\right)  -\sum_{1\leq j\leq2s}\varepsilon_{j}\left(
\left\langle \left.  \psi_{j}\right\vert \psi_{j}\right\rangle -1\right)
-\eta\left(  \sum_{1\leq j\leq s}n_{j}-1\right)  . \label{larE}%
\end{equation}
Some comments on the minimization problem eq. $\left(  \ref{min}\right)  $:

\begin{enumerate}
\item The number of electrons being described is implicit in the definition of
$\mathcal{E}\left(  \mathbf{n,\psi,\theta}\right)  $ and one could show this
dependency explicitly as
\begin{equation}
\mathcal{E}=\mathcal{E}\left(  2N,\mathbf{n,\psi,\theta}\right)  .
\end{equation}


\item The constraint eq. $\left(  \ref{thetacon}\right)  $ is always satisfied
as $\mathcal{E}$ is a periodic function of $\left\{  \left.  \theta
_{j}\right\vert 2\leq j\leq s\right\}  $ with period $2\pi.$

\item We show in section \ref{procedure} that the normalization constraints
eq. $\left(  \ref{norm}\right)  $ imply the orthonormality constraints
\begin{equation}
\left\langle \left\vert \psi_{j}\right\vert \psi_{k}\right\rangle =\delta
_{jk};\quad1\leq j\leq2s.
\end{equation}


\item In the optimization procedure we describe in section \ref{procedure} we
allow the constraint eq. $\left(  \ref{rep}\right)  $ to relax, however we
show in section \ref{Evariation} that this does effect the final solution.

\item The positivity constraints eq. $\left(  \ref{pos1}\right)  $ can be
handled by defining
\begin{equation}
n_{j}=\sin^{2}\alpha_{j};\quad0<\alpha_{j}\leq\pi;\quad1\leq j\leq s
\end{equation}
producing an unconstrained variation over $\left\{  \alpha_{j};\quad1\leq
j\leq s\right\}  $.
\end{enumerate}

A necessary condition for $\left(  \mathbf{n\left(  \mathbf{\alpha}%
_{\#}\right)  ,\psi}_{\#}\mathbf{,\theta}_{\#}\right)  $ to be a solution of
eq. $\left(  \ref{min}\right)  $ is that there exists Lagrange multipliers
$\eta,\left\{  \left.  \varepsilon_{j}\right\vert 1\leq j\leq2s\right\}  $
such that
\begin{equation}
\partial^{1}\mathcal{\tilde{E}}\left(  \mathbf{n}\left(  \mathbf{\alpha
}\right)  _{\#}\mathbf{,\psi}_{\#}\mathbf{,\theta}_{\#}\right)  =\mathbf{0,}%
\end{equation}
where
\begin{equation}
\partial^{1}\mathcal{\tilde{E}}\left(  \mathbf{n}\left(  \mathbf{\alpha}%
_{\#}\right)  \mathbf{,\psi}_{\#}\mathbf{,\theta}_{\#}\right)  =\left(
\partial_{\mathbf{n}}^{1}\mathcal{\tilde{E}}\left(  \mathbf{n}\left(
\mathbf{\alpha}_{\#}\right)  \mathbf{,\psi}_{\#}\mathbf{,\theta}_{\#}\right)
,\partial_{\mathbf{\psi}}^{1}\mathcal{\tilde{E}}\left(  \mathbf{n}\left(
\mathbf{\alpha}_{\#}\right)  \mathbf{,\psi}_{\#}\mathbf{,\theta}_{\#}\right)
,\partial_{\mathbf{\theta}}^{1}\mathcal{\tilde{E}}\left(  \mathbf{n}\left(
\mathbf{\alpha}_{\#}\right)  \mathbf{,\psi}_{\#}\mathbf{,\theta}_{\#}\right)
\right)  ,
\end{equation}
and
\begin{align}
\partial_{\mathbf{n}}^{1}\mathcal{\tilde{E}}\left(  \mathbf{n}\left(
\mathbf{\alpha}_{\#}\right)  \mathbf{,\psi}_{\#}\mathbf{,\theta}_{\#}\right)
_{j}  &  =\frac{\partial\mathcal{\tilde{E}}\left(  \mathbf{n}\left(
\mathbf{\alpha}_{\#}\right)  \mathbf{,\psi}_{\#}\mathbf{,\theta}_{\#}\right)
}{\partial n_{j}};\quad1\leq j\leq s\nonumber\\
\partial_{\mathbf{\psi}}^{1}\mathcal{\tilde{E}}\left(  \mathbf{n}\left(
\mathbf{\alpha}_{\#}\right)  \mathbf{,\psi}_{\#}\mathbf{,\theta}_{\#}\right)
_{j}  &  =\frac{\delta\mathcal{\tilde{E}}\left(  \mathbf{n}\left(
\mathbf{\alpha}_{\#}\right)  \mathbf{,\psi}_{\#}\mathbf{,\theta}_{\#}\right)
}{\delta\psi_{j}^{\ast}};\quad1\leq j\leq s\nonumber\\
\partial_{\mathbf{\psi}}^{1}\mathcal{\tilde{E}}\left(  \mathbf{n}\left(
\mathbf{\alpha}_{\#}\right)  \mathbf{,\psi}_{\#}\mathbf{,\theta}_{\#}\right)
_{j+s}  &  =\frac{\delta\mathcal{\tilde{E}}\left(  \mathbf{n}\left(
\mathbf{\alpha}_{\#}\right)  \mathbf{,\psi}_{\#}\mathbf{,\theta}_{\#}\right)
}{\delta\psi_{j}};\quad1\leq j\leq s\nonumber\\
\partial_{\mathbf{\theta}}^{1}\mathcal{\tilde{E}}\left(  \mathbf{n}\left(
\mathbf{\alpha}_{\#}\right)  \mathbf{,\psi}_{\#}\mathbf{,\theta}_{\#}\right)
_{j}  &  =\frac{\partial\mathcal{\tilde{E}}\left(  \mathbf{n}\left(
\mathbf{\alpha}_{\#}\right)  \mathbf{,\psi}_{\#}\mathbf{,\theta}_{\#}\right)
}{\partial\theta_{j}};\quad2\leq j\leq s
\end{align}
and
\begin{align}
\left\langle \left.  \psi_{\#j}\right\vert \psi_{\#k}\right\rangle  &
=\delta_{jk};\quad1\leq j,k\leq s\nonumber\\
\sum_{1\leq j\leq s}n_{j}\left(  \alpha_{\#j}\right)   &  =1.
\end{align}
As is well known this condition only characterizes constrained stationary
points of $\mathcal{E}$ and the Hessian matrix of second derivatives has to be
examined in order to determine if the point $\left(  \mathbf{n}\left(
\mathbf{\alpha}_{\#}\right)  \mathbf{,}\left\vert \mathbf{\psi}_{\#}%
\right\rangle \mathbf{,\theta}_{\#}\right)  $ does in fact minimize
$\mathcal{E}.$

\section{Generalized Fock Operator \label{Fock}}

\subsection{Functional Derivatives}

The energy expression eq. $\left(  \ref{energy}\right)  $ can be
differentiated wrt $\mathbf{\psi}^{\ast},$ $\mathbf{\psi}$ to produce the
functional derivatives (Appendix )
\begin{equation}
\left\{  \frac{\delta\mathcal{E}\left(  \mathbf{n,\psi,\theta}\right)
}{\delta\psi_{j}^{\ast}},\frac{\delta\mathcal{E}\left(  \mathbf{n,\psi,\theta
}\right)  }{\delta\psi_{j}};\quad1\leq j\leq s\right\}  ,
\end{equation}
which can be expressed as%
\begin{align}
\frac{\delta\mathcal{E}\left(  \mathbf{n,\psi,\theta}\right)  }{\delta\psi
_{j}^{\ast}}  &  =\frac{1}{S_{N}\left(  \mathbf{n}\right)  }\frac
{\delta\mathcal{H}\left(  \mathbf{n,\psi,\theta}\right)  }{\delta\psi
_{j}^{\ast}}\\
\frac{\delta\mathcal{E}\left(  \mathbf{n,\psi,\theta}\right)  }{\delta\psi
_{j}}  &  =\frac{1}{S_{N}\left(  \mathbf{n}\right)  }\frac{\delta
\mathcal{H}\left(  \mathbf{n,\psi,\theta}\right)  }{\delta\psi_{j}}%
\end{align}
in terms of the derivatives of the unnormalized energy $\mathcal{H}\left(
\mathbf{n,\psi,\theta}\right)  $ as $S_{N}\left(  \mathbf{n}\right)  $ does
not depend on $\left(  \mathbf{\psi,\theta}\right)  $ when $\left\{  \left.
\left\vert \psi_{j}\right\rangle \right\vert \quad1\leq j\leq2s\right\}  $ are
orthonormal. It is easy to show that
\begin{equation}
\frac{\delta\mathcal{H}\left(  \mathbf{n,\psi,\theta}\right)  }{\delta\psi
_{j}^{\ast}}=\left(  \frac{\delta\mathcal{H}\left(  \mathbf{n,\psi,\theta
}\right)  }{\delta\psi_{j}}\right)  ^{\ast}%
\end{equation}
so we only need consider $\left\{  \frac{\delta\mathcal{H}\left(
\mathbf{n,\psi,\theta}\right)  }{\delta\psi_{j}^{\ast}}\right\}  .$ These
derivatives can be expressed $\left(  \text{see Appendix\ref{variational}%
}\right)  $ in terms of a generalized self adjoint Fock operator $F\left(
\mathbf{n,\psi,\theta}\right)  $ as
\begin{equation}
\frac{\delta\mathcal{H}\left(  \mathbf{n,\psi,\theta}\right)  }{\delta\psi
_{j}^{\ast}}=F\left(  \mathbf{n,\psi,\theta}\right)  \left\vert \psi
_{j}\right\rangle ,
\end{equation}
that depends on the occupation numbers $\mathbf{n,}$ representative CGSO's
$\mathbf{\psi,}$of the class $\left[  D_{2N}^{1}\right]  $ and the internal
phases $\mathbf{\theta}$ and is given by
\begin{equation}
F\left(  \mathbf{n,\psi,\theta}\right)  =\sum_{1\leq j,m\leq2s}\left\{
N_{j}\left\{  h_{1mj}+A_{mj}-B_{mj}\right\}  +\left\{  C_{mj}-D_{mj}\right\}
+2\overset{}{\left\{  E_{mj}-F_{mj}\right\}  }\right\}  \left\vert \psi
_{m}\right\rangle \left\langle \psi_{j}\right\vert .
\end{equation}
The matrices $\left\{  \mathbb{A},\mathbb{B},\mathbb{C},\mathbb{D}%
,\mathbb{E},\mathbb{F}\right\}  $ are given by in terms of the matrix elements
$\left\{  \left.  \left(  \left.  \psi_{j}\right\vert h_{2}\psi_{k}\right)
_{lm}\right\vert 1\leq l,m\leq2s\right\}  $ of the operators $\left\{  \left.
\left(  \left.  \psi_{j}\right\vert h_{2}\psi_{k}\right)  \right\vert 1\leq
j,k\leq2s\right\}  ,$ that are defined by
\begin{equation}
\left(  \left.  \psi_{j}\right\vert h_{2}\psi_{k}\right)  \left(  \left\vert
\psi_{1}\right\rangle ,\cdots,\left\vert \psi_{2s}\right\rangle \right)
=\left(  \left\vert \psi_{1}\right\rangle ,\cdots,\left\vert \psi
_{2s}\right\rangle \right)  \mathbb{M}^{jk},
\end{equation}
i.e.
\begin{equation}
\left(  \left.  \psi_{j}\right\vert h_{2}\psi_{k}\right)  \left\vert \psi
_{m}\right\rangle =%
%TCIMACRO{\tsum \limits_{1\leq l\leq2s}}%
%BeginExpansion
{\textstyle\sum\limits_{1\leq l\leq2s}}
%EndExpansion
\left\vert \psi_{l}\right\rangle \mathbb{M}_{lm}^{jk};\quad1\leq m\leq2s,
\end{equation}
where
\begin{equation}
\left(  \mathbb{M}^{jk}\right)  _{lm}=\int\psi_{l}^{\ast}\left(
\mathbf{x}_{1}\right)  \psi_{j}^{\ast}\left(  \mathbf{x}_{2}\right)
h_{2}\left(  \mathbf{x}_{1}\mathbf{x}_{2};\mathbf{x}_{1}^{\prime}%
\mathbf{x}_{2}^{\prime}\right)  \psi_{m}\left(  \mathbf{x}_{1}^{\prime
}\right)  \psi_{k}\left(  \mathbf{x}_{2}^{\prime}\right)  d\mathbf{x}%
_{1}d\mathbf{x}_{2}d\mathbf{x}_{1}^{\prime}d\mathbf{x}_{2}^{\prime}.
\end{equation}
The matrix elements of $\left\{  \mathbb{A},\mathbb{B},\mathbb{C}%
,\mathbb{D},\mathbb{E},\mathbb{F}\right\}  $ for $1\leq m\leq2s$ are given by
\begin{equation}
A_{mj}=\left\{
\begin{array}
[c]{c}%
\left(  \left.  \psi_{j+s}\right\vert h_{2}\psi_{j+s}\right)  _{mj};1\leq
j\leq s\\
\left(  \left.  \psi_{j-s}\right\vert h_{2}\psi_{j-s}\right)  _{mj};s+1\leq
j\leq s
\end{array}
\right.  ,
\end{equation}%
\begin{equation}
B_{mj}=\left\{
\begin{array}
[c]{c}%
\left(  \left.  \psi_{j+s}\right\vert h_{2}\psi_{j}\right)  _{mj+s};1\leq
j\leq s\\
\left(  \left.  \psi_{j-s}\right\vert h_{2}\psi_{j}\right)  _{mj-s};s+1\leq
j\leq s
\end{array}
\right.  ,
\end{equation}
noting that $n_{j+s}=n_{j}$ and $\theta_{j+s}=\theta_{j}$ for $1\leq j\leq s$
we have the phase dependent $2$-electron terms
\begin{equation}
C_{mj}=\left\{
\begin{array}
[c]{c}%
%TCIMACRO{\tsum \limits_{1\leq k\leq s}}%
%BeginExpansion
{\textstyle\sum\limits_{1\leq k\leq s}}
%EndExpansion
n_{j}^{\frac{1}{2}}n_{k}^{\frac{1}{2}}S_{N-1}\left(  \jmath k\right)
e^{i\left(  \theta_{j}-\theta_{k}\right)  }\left(  \left.  \psi_{j+s}%
\right\vert h_{2}\psi_{k+s}\right)  _{mk};1\leq j\leq s\\%
%TCIMACRO{\tsum \limits_{1\leq k\leq s}}%
%BeginExpansion
{\textstyle\sum\limits_{1\leq k\leq s}}
%EndExpansion
n_{j}^{\frac{1}{2}}n_{k}^{\frac{1}{2}}S_{N-1}\left(  \jmath k\right)
e^{i\left(  \theta_{j}-\theta_{k}\right)  }\left(  \left.  \psi_{j-s}%
\right\vert h_{2}\psi_{k}\right)  _{mk+s};s+1\leq j\leq2s
\end{array}
\right.  ,
\end{equation}%
\begin{equation}
D_{mj}=\left\{
\begin{array}
[c]{c}%
%TCIMACRO{\tsum \limits_{1\leq k\leq s}}%
%BeginExpansion
{\textstyle\sum\limits_{1\leq k\leq s}}
%EndExpansion
n_{j}^{\frac{1}{2}}n_{k}^{\frac{1}{2}}S_{N-1}\left(  \jmath k\right)
e^{i\left(  \theta_{j}-\theta_{k}\right)  }\left(  \left.  \psi_{j+s}%
\right\vert h_{2}\psi_{k}\right)  _{mk+s};1\leq j\leq s\\%
%TCIMACRO{\tsum \limits_{1\leq k\leq s}}%
%BeginExpansion
{\textstyle\sum\limits_{1\leq k\leq s}}
%EndExpansion
n_{j}^{\frac{1}{2}}n_{k}^{\frac{1}{2}}S_{N-1}\left(  \jmath k\right)
e^{i\left(  \theta_{j}-\theta_{k}\right)  }\left(  \left.  \psi_{j-s}%
\right\vert h_{2}\psi_{k}\right)  _{mk+s};s+1\leq j\leq2s
\end{array}
\right.  ,
\end{equation}
the phase independent $2$-electron terms
\begin{equation}
F_{mj}=\sum_{1\leq k\leq2s}n_{j}n_{k}S_{N-2}\left(  \jmath k\right)  \left(
\left.  \psi_{k}\right\vert h_{2}\psi_{j}\right)  _{mk};1\leq j\leq2s,
\end{equation}
and%
\begin{equation}
E_{mj}=\sum_{1\leq k\leq2s}n_{j}n_{k}S_{N-2}\left(  \jmath k\right)  \left(
\left.  \psi_{k}\right\vert h_{2}\psi_{k}\right)  _{mj};1\leq j\leq2s.
\end{equation}


\subsection{Orthonormality Constraints}

A necessary condition for $\mathbf{\psi}_{\ddag}$ $=\mathbf{\psi}_{\ddag
}\left(  \mathbf{n,\theta}\right)  $ to be a orthonormal constrained
stationary point of $\mathcal{E}\left(  \mathbf{n,\psi,\theta}\right)  $ for a
fixed value of $\left(  \mathbf{n,\theta}\right)  $ is given by
\begin{subequations}
\begin{align}
\partial_{\mathbf{\psi}}^{1}\mathcal{\tilde{H}}\left(  \mathbf{n,\mathbf{\psi
}_{\ddag},\theta}\right)   &  =\mathbf{0}\label{nec1}\\
\left\langle \left.  \psi_{\ddag j}\right\vert \psi_{\ddag k}\right\rangle  &
=\delta_{jk};1\leq j,k\leq2s,
\end{align}
where the Lagrangian $\mathcal{\tilde{H}}\left(  \mathbf{n,\mathbf{\psi
},\theta}\right)  $ is defined as
\end{subequations}
\begin{equation}
\mathcal{\tilde{H}}\left(  \mathbf{n,\mathbf{\psi},\theta}\right)
=\mathcal{H}\left(  \mathbf{n,\mathbf{\psi},\theta}\right)  -\sum_{1\leq
j\leq2s}\varepsilon_{j}\left\{  \left\langle \left.  \psi_{j}\right\vert
\psi_{j}\right\rangle -1\right\}  , \label{lagH}%
\end{equation}
as
\begin{equation}
\partial_{\mathbf{\psi}}^{1}\mathcal{\tilde{H}}\left(  \mathbf{n,\mathbf{\psi
},\theta}\right)  =\partial_{\mathbf{\psi}}^{1}\mathcal{\tilde{E}}\left(
\mathbf{n,\mathbf{\psi},\theta}\right)
\end{equation}
for orthonormality constrained $\mathbf{\mathbf{\psi.}}$ The condition eq.
$\left(  \ref{nec1}\right)  $ can be expressed in terms of the generalized
Fock operator as
\begin{equation}
F\left(  \mathbf{n,\mathbf{\psi},\theta}\right)  \left\vert \psi
_{j}\right\rangle =\varepsilon_{j}\left\vert \psi_{j}\right\rangle ;1\leq
j\leq2s
\end{equation}
As $F\left(  \mathbf{\psi,n}\right)  =F\left(  \mathbf{\psi,n}\right)
^{\dagger}$ this implies that one can find solutions $\left\{  \left.
\left\vert \psi_{m}\right\rangle \right\vert 1\leq m\leq2s\right\}  $ for
$\left\{  \left.  \varepsilon_{m}\right\vert 1\leq m\leq2s\right\}  $ such
that
\begin{equation}
\left\langle \left.  \psi_{j}\right\vert \psi_{k}\right\rangle =\delta
_{jk},;1\leq j,k\leq2s.
\end{equation}
Thus one can always find stationary points that satisfy the \emph{normality}
constraints
\begin{equation}
\left\langle \left.  \psi_{j}\right\vert \psi_{j}\right\rangle =1;1\leq
j\leq2s
\end{equation}
that also satisfy the orthonormality constraints eq. $\left(  \ref{norm}%
\right)  .$

The Lagrangian eq. $\left(  \ref{lagH}\right)  $ can also be used in the
IPS\ case to produce the standard Hartree-Fock equations for $2N$ -electrons
instead of the traditional Lagrangian that explicitly holds the orthonormality
constraints between occupied orbitals
\begin{equation}
\mathcal{\tilde{H}}\left(  \mathbf{n}_{IP}\mathbf{,\mathbf{\psi},0}\right)
=\mathcal{H}\left(  \mathbf{n}_{IP}\mathbf{,\mathbf{\psi},0}\right)
-\sum_{1\leq j,k\leq2N}\varepsilon_{jk}\left\{  \left\langle \left.  \psi
_{j}\right\vert \psi_{k}\right\rangle -1\right\}  . \label{tradlag}%
\end{equation}
\emph{However }using a Lagrangian that explicitly holds the orthonormality
constraints in the AGP\ case \emph{does not} lead to equations that impose
conditions on $\mathbf{\mathbf{\psi}}$.

\section{Energy Variation Procedure \label{procedure}}

\section{Discussion\label{Discussion}}

\section{Acknowledgements}

\appendix

\section{One Electron Basis\label{basis}}

\section{Example of Linear Fit\ref{linfit}}

consider
\begin{align}
\mathbf{N}_{1}  &  =\left(
\begin{array}
[c]{c}%
1\\
.9\\
.8\\
.1\\
.2\\
0
\end{array}
\right)  =\left(
\begin{array}
[c]{c}%
1\\
1\\
1\\
0\\
0\\
0
\end{array}
\right)  +\left(
\begin{array}
[c]{c}%
0\\
-.1\\
-.2\\
.1\\
.2\\
0
\end{array}
\right)  =\mathbf{N}_{0}+\delta\mathbf{N}_{1}\\
\mathbf{n}_{1}  &  =\mathcal{N}^{-1}\left(
\begin{array}
[c]{c}%
1\\
.9\\
.8\\
0\\
.2\\
.1
\end{array}
\right)  =\left(
\begin{array}
[c]{c}%
\frac{1}{3}\\
\frac{1}{3}\\
\frac{1}{3}\\
0\\
0\\
0
\end{array}
\right)  +\left[  d\mathcal{N}\left(
\begin{array}
[c]{c}%
\frac{1}{3}\\
\frac{1}{3}\\
\frac{1}{3}\\
0\\
0\\
0
\end{array}
\right)  \right]  ^{-1}\left(
\begin{array}
[c]{c}%
0\\
-.1\\
-.2\\
.1\\
.2\\
0
\end{array}
\right) \\
\mathbf{n}_{2}  &  =\mathcal{N}^{-1}\left(
\begin{array}
[c]{c}%
.9\\
.8\\
.8\\
.1\\
.2\\
.2
\end{array}
\right)  =\left(
\begin{array}
[c]{c}%
\frac{1}{3}\\
\frac{1}{3}\\
\frac{1}{3}\\
0\\
0\\
0
\end{array}
\right)  +\left[  d\mathcal{N}\left(
\begin{array}
[c]{c}%
\frac{1}{3}\\
\frac{1}{3}\\
\frac{1}{3}\\
0\\
0\\
0
\end{array}
\right)  \right]  ^{-1}\left(
\begin{array}
[c]{c}%
0\\
-.1\\
-.2\\
.1\\
.2\\
0
\end{array}
\right) \\
&  +\left[  d\mathcal{N}\left(  \left(
\begin{array}
[c]{c}%
\frac{1}{3}\\
\frac{1}{3}\\
\frac{1}{3}\\
0\\
0\\
0
\end{array}
\right)  +\left[  d\mathcal{N}\left(
\begin{array}
[c]{c}%
\frac{1}{3}\\
\frac{1}{3}\\
\frac{1}{3}\\
0\\
0\\
0
\end{array}
\right)  \right]  ^{-1}\left(
\begin{array}
[c]{c}%
0\\
-.1\\
-.2\\
.1\\
.2\\
0
\end{array}
\right)  \right)  \right]  ^{-1}\left(
\begin{array}
[c]{c}%
-.1\\
-.1\\
0\\
.1\\
0\\
.1
\end{array}
\right) \\
&  \equiv\mathbf{n}_{0}+\left[  d\mathcal{N}\left(  \mathbf{n}_{0}\right)
\right]  ^{-1}\left(  \delta\mathbf{N}_{1}\right)  +\left[  d\mathcal{N}%
\left(  \mathbf{n}_{0}+\left[  d\mathcal{N}\left(  \mathbf{n}_{0}\right)
\right]  ^{-1}\left(  \delta\mathbf{N}_{1}\right)  \right)  \right]
^{-1}\left(  \left(  \delta\mathbf{N}_{2}\right)  \right)
\end{align}


\section{Variational Equations\label{variational}}

\subsection{Derivatives}

The $2N$-AGP\ energy functional in terms of the CGSO's $\left\{  \left.
\psi_{j}\right\vert 1\leq j\leq2s\right\}  $, the geminal occupation numbers
$\left\{  \left.  n_{j}\right\vert 1\leq j\leq s\right\}  $ and the phases
$\left\{  \left.  \theta_{j}\right\vert 1\leq j\leq s\right\}  $ can be
expressed as the quotient%
\begin{align}
\mathcal{E}\left(  \mathbf{n},\mathbf{\psi},\mathbf{\theta}\right)   &
=\frac{\mathcal{H}\left(  \mathbf{n},\mathbf{\psi},\mathbf{\theta}\right)
}{S_{N}\left(  \mathbf{n}\right)  }=\frac{1}{S_{N}\left(  \mathbf{n}\right)
}\left\{  \sum_{1\leq j\leq s}N_{j}\left\{  \left\langle \psi_{j}\left\vert
h_{1}\psi_{j}\right.  \right\rangle +\left\langle \psi_{j+s}\left\vert
h_{1}\psi_{j+s}\right.  \right\rangle +\left\langle \psi_{j}\psi_{j+s}%
|h_{2}|\psi_{j}\psi_{j+s}\right\rangle \right\}  \right. \nonumber\\
&  +\left.  \sum_{1\leq j,k\leq s}n_{j}^{\frac{1}{2}}n_{k}^{\frac{1}{2}%
}e^{i\left(  \theta_{j}-\theta_{k}\right)  }S_{N-1}\left(  \jmath k\right)
\left\langle \psi_{j}\psi_{j+s}|h_{2}|\psi_{k}\psi_{k+s}\right\rangle
+\sum_{1\leq j,k\leq s}n_{j}n_{k}S_{N-2}\left(  \jmath k\right)  \left\langle
\psi_{j}\psi_{k}|||\psi_{j}\psi_{k}\right\rangle \right\}  .
\end{align}
The first order variation $\delta^{1}\mathcal{E}\left(  \mathbf{n}%
,\mathbf{\psi},\mathbf{\theta}\right)  $ of the energy is then%
\begin{equation}
\delta^{1}\mathcal{E}\left(  \mathbf{n},\mathbf{\psi},\mathbf{\theta}\right)
=\left(
\begin{array}
[c]{cccc}%
\partial_{\mathbf{n}}^{1}\mathcal{E}\left(  \mathbf{n},\mathbf{\psi
},\mathbf{\theta}\right)  & \partial_{\mathbf{\psi}^{\ast}}^{1}\mathcal{E}%
\left(  \mathbf{n},\mathbf{\psi},\mathbf{\theta}\right)  & \partial
_{\mathbf{\psi}}^{1}\mathcal{E}\left(  \mathbf{n},\mathbf{\psi},\mathbf{\theta
}\right)  & \partial_{\mathbf{\theta}}^{1}\mathcal{E}\left(  \mathbf{n}%
,\mathbf{\psi},\mathbf{\theta}\right)
\end{array}
\right)  \left(
\begin{array}
[c]{c}%
\delta\mathbf{n}\\
\delta\mathbf{\psi}^{\ast}\\
\delta\mathbf{\psi}\\
\delta\mathbf{\theta}%
\end{array}
\right)  ,
\end{equation}
where
\begin{subequations}
\begin{align}
\partial_{\mathbf{n}}^{1}\mathcal{E}\left(  \mathbf{n},\mathbf{\psi
},\mathbf{\theta}\right)   &  =\frac{1}{S_{N}\left(  \mathbf{n}\right)
}\left\{  \partial_{\mathbf{n}}^{1}\mathcal{H}\left(  \mathbf{n},\mathbf{\psi
},\mathbf{\theta}\right)  -\mathcal{E}\left(  \mathbf{n},\mathbf{\psi
},\mathbf{\theta}\right)  \partial_{\mathbf{n}}^{1}S_{N}\left(  \mathbf{n}%
\right)  \right\} \\
\partial_{\mathbf{\psi}^{\ast}}^{1}\mathcal{E}\left(  \mathbf{n},\mathbf{\psi
},\mathbf{\theta}\right)   &  =\frac{1}{S_{N}\left(  \mathbf{n}\right)
}\left\{  \partial_{\mathbf{\psi}^{\ast}}^{1}\mathcal{H}\left(  \mathbf{n}%
,\mathbf{\psi},\mathbf{\theta}\right)  -\mathcal{E}\left(  \mathbf{n}%
,\mathbf{\psi},\mathbf{\theta}\right)  \partial_{\mathbf{\psi}^{\ast}}%
^{1}S_{N}\left(  \mathbf{n}\right)  \right\}  =\left[  \partial_{\mathbf{\psi
}}^{1}\mathcal{E}\left(  \mathbf{n},\mathbf{\psi},\mathbf{\theta}\right)
\right]  ^{\ast}\\
\partial_{\mathbf{\theta}}^{1}\mathcal{E}\left(  \mathbf{n},\mathbf{\psi
},\mathbf{\theta}\right)   &  =\frac{1}{S_{N}\left(  \mathbf{n}\right)
}\left\{  \partial_{\mathbf{\theta}}^{1}\mathcal{H}\left(  \mathbf{n}%
,\mathbf{\psi},\mathbf{\theta}\right)  -\mathcal{E}\left(  \mathbf{n}%
,\mathbf{\psi},\mathbf{\theta}\right)  \partial_{\mathbf{\theta}}^{1}%
S_{N}\left(  \mathbf{n}\right)  \right\}
\end{align}
and
\end{subequations}
\begin{subequations}
\begin{align}
\partial_{\mathbf{\psi}^{\ast}}^{1}\mathcal{E}\left(  \mathbf{n},\mathbf{\psi
},\mathbf{\theta}\right)  \cdot\delta\mathbf{\psi}^{\ast}  &  =\sum_{1\leq
j\leq2s}\partial_{\mathbf{\psi}^{\ast}}^{1}\mathcal{E}\left(  \mathbf{n}%
,\mathbf{\psi},\mathbf{\theta}\right)  \left(  \delta\psi_{j}^{\ast}\right)
=\sum_{1\leq j\leq2s}\int\frac{\delta\mathcal{E}\left(  \mathbf{n}%
,\mathbf{\psi},\mathbf{\theta}\right)  }{\delta\psi_{j}^{\ast}\left(
\mathbf{r,}\zeta\right)  }\delta\psi_{j}^{\ast}\left(  \mathbf{r,}%
\zeta\right)  d\mathbf{r}d\zeta\label{funcder}\\
\partial_{\mathbf{n}}^{1}\mathcal{E}\left(  \mathbf{n},\mathbf{\psi
},\mathbf{\theta}\right)  \cdot\delta\mathbf{n}  &  =\sum_{1\leq j\leq s}%
\frac{\partial\mathcal{E}\left(  \mathbf{n},\mathbf{\psi},\mathbf{\theta
}\right)  }{\partial n_{j}}\delta n_{j}.
\end{align}
If the variations $\left\{  \delta\mathbf{\psi}^{\ast},\delta\mathbf{\psi
}\right\}  $ maintain the orthonormality of the CGSO's i.e.
\end{subequations}
\begin{equation}
\left\langle \left.  \psi_{j}\right\vert \psi_{k}\right\rangle =\delta_{jk}%
\end{equation}
then
\begin{align}
\partial_{\mathbf{\psi}^{\ast}}^{1}\mathcal{E}\left(  \mathbf{n},\mathbf{\psi
},\mathbf{\theta}\right)   &  =\frac{1}{S_{N}\left(  \mathbf{n}\right)
}\partial_{\mathbf{\psi}^{\ast}}^{1}\mathcal{H}\left(  \mathbf{n}%
,\mathbf{\psi},\mathbf{\theta}\right) \\
\partial_{\mathbf{\theta}}^{1}\mathcal{E}\left(  \mathbf{n},\mathbf{\psi
},\mathbf{\theta}\right)   &  =\frac{1}{S_{N}\left(  \mathbf{n}\right)
}\partial_{\mathbf{\theta}}^{1}\mathcal{H}\left(  \mathbf{n},\mathbf{\psi
},\mathbf{\theta}\right)
\end{align}
as $\partial_{\mathbf{\psi}^{\ast}}^{1}S_{N}\left(  \mathbf{n}\right)
=\partial_{\mathbf{\psi}}^{1}S_{N}\left(  \mathbf{n}\right)  =\mathbf{0}$.

\subsection{Constraints on the Geminal Occupation Numbers}

In the variation of $\mathbf{n}$ one needs to hold the constraints
\begin{align}
&  \sum_{1\leq j\leq s}n_{j}=1\label{numeq}\\
&  0\leq n_{j}\leq1;\quad1\leq j\leq s \label{numpos}%
\end{align}
This constraint can be held by introducing a Lagrange multiplier $\eta$ in the
Lagrangian
\begin{equation}
\mathcal{\tilde{E}}\left(  \mathbf{n},\mathbf{\psi},\mathbf{\theta}\right)
=\mathcal{E}\left(  \mathbf{n},\mathbf{\psi},\mathbf{\theta}\right)
-\sum_{1\leq j\leq2s}\varepsilon_{j}\left(  \left\langle \left.  \psi
_{j}\right\vert \psi_{j}\right\rangle -1\right)  +\eta\left(  \sum_{1\leq
j\leq s}n_{j}-1\right)
\end{equation}
producing
\begin{equation}
\partial_{n_{j}}^{1}\mathcal{\tilde{E}}\left(  \mathbf{n},\mathbf{\psi
},\mathbf{\theta}\right)  =\frac{1}{S_{N}\left(  \mathbf{n}\right)  }\left\{
\partial_{n_{j}}^{1}\mathcal{H}\left(  \mathbf{n},\mathbf{\psi},\mathbf{\theta
}\right)  -\mathcal{E}\left(  \mathbf{n},\mathbf{\psi},\mathbf{\theta}\right)
\partial_{n_{j}}^{1}S_{N}\left(  \mathbf{n}\right)  +\eta S_{N}\left(
\mathbf{n}\right)  \right\}
\end{equation}
which leads to
\begin{equation}
\partial_{n_{j}}^{1}\mathcal{\tilde{E}}\left(  \mathbf{n},\mathbf{\psi
},\mathbf{\theta}\right)  =\frac{1}{S_{N}\left(  \mathbf{n}\right)  }\left\{
%TCIMACRO{\tsum \limits_{k}}%
%BeginExpansion
{\textstyle\sum\limits_{k}}
%EndExpansion
\partial_{n_{j}n_{k}}^{2}\mathcal{H}\left(  \mathbf{n},\mathbf{\psi
},\mathbf{\theta}\right)  n_{k}-\mathcal{E}\left(  \mathbf{n},\mathbf{\psi
},\mathbf{\theta}\right)
%TCIMACRO{\tsum \limits_{k}}%
%BeginExpansion
{\textstyle\sum\limits_{k}}
%EndExpansion
\partial_{n_{j}n_{k}}^{2}S_{N}\left(  \mathbf{n}\right)  \right\}  n_{k}+\eta%
%TCIMACRO{\tsum \limits_{j,k}}%
%BeginExpansion
{\textstyle\sum\limits_{j,k}}
%EndExpansion
\partial_{n_{j}n_{k}}^{2}S_{N}\left(  \mathbf{n}\right)  n_{j}n_{k}%
\end{equation}
and an eigenvalue problem *********

The positivity constraint is handled by introducing
\begin{equation}
n_{j}=\sin^{2}\alpha_{j}%
\end{equation}
and noting that
\begin{equation}
\frac{\partial f\left(  \mathbf{n}\left(  \mathbf{\alpha}\right)  \right)
}{\partial\alpha_{j}}=\sum\limits_{k}\frac{\partial f\left(  \mathbf{n}\left(
\mathbf{\alpha}\right)  \right)  }{\partial n_{k}}\frac{\partial n_{k}\left(
\mathbf{\alpha}\right)  }{\partial\alpha_{j}}=2\sum\limits_{k}\frac{\partial
f\left(  \mathbf{n}\left(  \mathbf{\alpha}\right)  \right)  }{\partial n_{k}%
}\sin\alpha_{j}\cos\alpha_{j}\delta_{jk}=\frac{\partial f\left(
\mathbf{n}\left(  \mathbf{\alpha}\right)  \right)  }{\partial n_{j}}%
\sin2\alpha_{j}.
\end{equation}


\subsection{Derivation of the Fock Operator}

In order to obtain derivative of terms in the energy expression it helpful to
utilize matrix elements wrt non antisymmetric states
\begin{equation}
\left\vert \psi_{j}\psi_{k}\right)  =\left\vert \psi_{j}\right)
\otimes\left\vert \psi_{k}\right)  ,
\end{equation}
whose amplitudes are
\begin{equation}
\left(  \left.  \mathbf{r}_{1}\mathbf{,\xi}_{1}\mathbf{r}_{2}\mathbf{,\xi}%
_{2}\right\vert \psi_{j}\psi_{k}\right)  =\psi_{j}\left(  \mathbf{r}%
_{1}\mathbf{,\xi}_{1}\right)  \psi_{k}\left(  \mathbf{r}_{2}\mathbf{,\xi}%
_{2}\right)  ,
\end{equation}
where
\begin{equation}
\psi_{j}\left(  \mathbf{r}_{1}\mathbf{,\xi}_{1}\right)  =\left(  \left.
\mathbf{r}_{1}\mathbf{,\xi}_{1}\right\vert \psi_{j}\right)  =\left\langle
\left.  \mathbf{r}_{1}\mathbf{,\xi}_{1}\right\vert \psi_{j}\right\rangle .
\end{equation}
In terms of these states the unnormalized $2N$-electron AGP\ energy becomes
\begin{align}
&  \mathcal{H}\left(  \mathbf{n},\mathbf{\psi},\mathbf{\theta}\right)
=\sum_{1\leq j\leq s}N_{j}\left\{  \left(  \left.  \psi_{j}\right\vert
h_{1}\psi_{j}\right)  +\left(  \left.  \psi_{j+s}\right\vert h_{1}\psi
_{j+s}\right)  +\left(  \left.  \psi_{j}\psi_{j+s}\right\vert h_{2}\psi
_{j}\psi_{j+s}\right)  -\left(  \left.  \psi_{j}\psi_{j+s}\right\vert
h_{2}\psi_{j+s}\psi_{j}\right)  \right\} \nonumber\\
&  +\sum_{1\leq j,k\leq s}\left\{  n_{j}^{\frac{1}{2}}n_{k}^{\frac{1}{2}%
}e^{i\left(  \theta_{j}-\theta_{k}\right)  }S_{N-1}\left(  \jmath k\right)
\left\{  \left(  \left.  \psi_{j}\psi_{j+s}\right\vert h_{2}\psi_{k}\psi
_{k+s}\right)  -\left(  \left.  \psi_{j}\psi_{j+s}\right\vert h_{2}\psi
_{k+s}\psi_{k}\right)  \right\}  \right. \nonumber\\
&  +n_{j}n_{k}S_{N-2}\left(  \jmath k\right)  \left\{  \overset{}{}\left(
\left.  \psi_{j}\psi_{k}\right\vert h_{2}\psi_{j}\psi_{k}\right)  -\left(
\left.  \psi_{j}\psi_{k}\right\vert h_{2}\psi_{k}\psi_{j}\right)  \right. \\
&  +\left(  \left.  \psi_{j}\psi_{k+s}\right\vert h_{2}\psi_{j}\psi
_{k+s}\right)  -\left(  \left.  \psi_{j}\psi_{k+s}\right\vert h_{2}\psi
_{k+s}\psi_{j}\right)  +\left(  \left.  \psi_{j+s}\psi_{k}\right\vert
h_{2}\psi_{j+s}\psi_{k}\right)  -\left(  \left.  \psi_{j+s}\psi_{k}\right\vert
h_{2}\psi_{k}\psi_{j+s}\right) \nonumber\\
&  \qquad\qquad\qquad\qquad\qquad\qquad\qquad\left.  \left.  +\left(  \left.
\psi_{j+s}\psi_{k+s}\right\vert h_{2}\psi_{j+s}\psi_{k+s}\right)  -\left(
\left.  \psi_{j+s}\psi_{k+s}\right\vert h_{2}\psi_{k+s}\psi_{j+s}\right)
\right\}  \overset{}{}\right\}  , \label{unnormE}%
\end{align}
where the form of $h_{1}$ and $h_{2}$ can be seen from eqs. $\left(
\ref{h1}\right)  $ and $\left(  \ref{h2}\right)  .$ In order to find the
derivatives of eq. $\left(  \ref{unnormE}\right)  $ one observes that
\begin{equation}
\frac{\delta}{\delta\psi_{\mu}^{\ast}}\left(  \left.  \psi_{a}\right\vert
h_{1}\psi_{b}\right)  =h_{1}\left\vert \psi_{b}\right\rangle \delta_{a\mu
}=\frame{$\left\{  h_1\left\vert \psi_b\right\rangle \left\langle
\psi_a\right\vert \right\}  \left\vert \psi_\mu\right\rangle $}=\left\{
\sum_{1\leq l,m\leq2s}h_{1lm}\left\vert \psi_{l}\right\rangle \left\langle
\psi_{m}\right\vert \right\}  \left\vert \psi_{b}\right\rangle \delta_{a\mu}%
\end{equation}
and
\begin{align}
&  \frac{\delta}{\delta\psi_{\mu}^{\ast}}\left(  \left.  \psi_{a}\psi
_{b}\right\vert h_{2}\psi_{c}\psi_{d}\right)  =\left(  \left.  \psi
_{b}\right\vert h_{2}\psi_{d}\right)  \left\vert \psi_{c}\right\rangle
\delta_{a\mu}+\left(  \left.  \psi_{a}\right\vert h_{2}\psi_{c}\right)
\left\vert \psi_{d}\right\rangle \delta_{b\mu}\nonumber\\
&  =\frame{$\left\{  \left(  \left.  \psi_b\right\vert h_2\psi_d\right)
\left\vert \psi_c\right\rangle \left\langle \psi_a\right\vert +\left(  \left.
\psi_a\right\vert h_2\psi_c\right)  \left\vert \psi_d\right\rangle
\left\langle \psi_b\right\vert \right\}  \left\vert \psi_\mu\right\rangle
$}\nonumber\\
&  =\sum_{1\leq l,m\leq2s}\left\{  \overset{}{\left(  \left.  \psi
_{b}\right\vert h_{2}\psi_{d}\right)  _{lm}\left\vert \psi_{l}\right\rangle
\left\langle \psi_{m}\right\vert }\right\}  \left\vert \psi_{c}\right\rangle
\delta_{a\mu}+\sum_{1\leq l,m\leq2s}\left\{  \overset{}{\left(  \left.
\psi_{a}\right\vert h_{2}\psi_{c}\right)  _{lm}\left\langle \psi
_{m}\right\vert }\right\}  \left\vert \psi_{d}\right\rangle \delta_{b\mu},
\end{align}
where
\begin{equation}
\left(  \left.  \psi_{b}\right\vert h_{2}\psi_{d}\right)  _{lm}=\int\psi
_{l}^{\ast}\left(  \mathbf{x}_{1}\right)  \psi_{b}^{\ast}\left(
\mathbf{x}_{2}\right)  h_{2}\left(  \mathbf{x}_{1}\mathbf{x}_{2}%
;\mathbf{x}_{1}^{\prime}\mathbf{x}_{2}^{\prime}\right)  \psi_{m}\left(
\mathbf{x}_{1}^{\prime}\right)  \psi_{d}\left(  \mathbf{x}_{2}^{\prime
}\right)  d\mathbf{x}_{1}d\mathbf{x}_{2}d\mathbf{x}_{1}^{\prime}%
d\mathbf{x}_{2}^{\prime}.
\end{equation}
Notice that the terms enclosed in boxes are \emph{products of operators}. If
we define the operators
\begin{align}
\left(  \left.  \psi_{j}\right\vert h_{2}\psi_{k}\right)  \left\vert \psi
_{l}\right\rangle  &  =\sum_{1\leq p,q\leq s}\left(  \left.  \psi
_{j}\right\vert h_{2}\psi_{k}\right)  _{pq}\left\vert \psi_{p}\right\rangle
\left\langle \left.  \psi_{q}\right\vert \psi_{l}\right\rangle =\left(
\left.  \cdot\psi_{j}\right\vert h_{2}\psi_{l}\psi_{k}\right) \\
\left(  \left.  \psi_{j}\right\vert h_{2}\psi_{k}\right)   &  =\sum_{1\leq
p,q\leq s}\left(  \left.  \psi_{j}\right\vert h_{2}\psi_{k}\right)
_{pq}\left\vert \psi_{p}\right\rangle \left\langle \psi_{q}\right\vert ,
\end{align}
where
\begin{align}
&  \left(  \left.  \psi_{j}\right\vert h_{2}\psi_{k}\right)  \left(
\mathbf{r}_{1},\mathbf{\xi}_{1}\right) \nonumber\\
&  =\int\psi_{j}^{\ast}\left(  \mathbf{r}_{2}\mathbf{,\xi}_{2}\right)
h_{2}\left(  \mathbf{r}_{1}\mathbf{,\xi}_{1};\mathbf{r}_{2}\mathbf{,\xi}%
_{2}\right)  \delta\left(  \mathbf{r}_{1}-\mathbf{r}_{1}^{\prime}\right)
\delta\left(  \mathbf{r}_{2}-\mathbf{r}_{2}^{\prime}\right)  \delta\left(
\mathbf{\xi}_{1}-\mathbf{\xi}_{1}^{\prime}\right)  \delta\left(  \mathbf{\xi
}_{2}-\mathbf{\xi}_{2}^{\prime}\right)  \psi_{k}\left(  \mathbf{r}%
_{2}\mathbf{,\xi}_{2}\right)  d\mathbf{r}_{2}d\mathbf{\xi}_{2}d\mathbf{r}%
_{2}^{\prime}d\mathbf{\xi}_{2}^{\prime}\nonumber\\
&  \left(  \left.  \psi_{j}\right\vert h_{2}\psi_{k}\right)  _{pq}\nonumber\\
&  =\int\psi_{p}^{\ast}\left(  \mathbf{r}_{1}\mathbf{,\xi}_{1}\right)
\psi_{j}^{\ast}\left(  \mathbf{r}_{2}\mathbf{,\xi}_{2}\right)  h_{2}\left(
\mathbf{r}_{1}\mathbf{,\xi}_{1};\mathbf{r}_{2}\mathbf{,\xi}_{2}\right)
\delta\left(  \mathbf{r}_{1}-\mathbf{r}_{1}^{\prime}\right)  \delta\left(
\mathbf{r}_{2}-\mathbf{r}_{2}^{\prime}\right)  \delta\left(  \mathbf{\xi}%
_{1}-\mathbf{\xi}_{1}^{\prime}\right)  \delta\left(  \mathbf{\xi}%
_{2}-\mathbf{\xi}_{2}^{\prime}\right) \nonumber\\
&  \qquad\qquad\psi_{k}\left(  \mathbf{r}_{2}^{\prime}\mathbf{,\xi}%
_{2}^{\prime}\right)  \psi_{q}\left(  \mathbf{r}_{1}^{\prime}\mathbf{,\xi}%
_{1}^{\prime}\right)  d\mathbf{r}_{2}d\mathbf{\xi}_{2}d\mathbf{r}_{2}^{\prime
}d\mathbf{\xi}_{2}^{\prime}%
\end{align}
The partial functional derivatives of $\mathcal{H}\left(  \mathbf{\psi
,n}\right)  $ wrt CGSO's $\mathbf{\psi}^{\ast}$ is then given by%
\begin{align}
&  \frac{\delta\mathcal{H}\left(  \mathbf{\psi,n}\right)  }{\delta\psi_{\mu
}^{\ast}}=\left\{  \sum_{1\leq j\leq s}N_{j}\left\{  h_{1}\left\{  \left\vert
\psi_{j}\right\rangle \left\langle \psi_{j}\right\vert +\left\vert \psi
_{j+s}\right\rangle \left\langle \psi_{j+s}\right\vert \right\}  +\left(
\left.  \psi_{j+s}\right\vert h_{2}\psi_{j+s}\right)  \left\vert \psi
_{j}\right\rangle \left\langle \psi_{j}\right\vert \right.  \right.
\nonumber\\
&  \qquad\left.  +\left(  \left.  \psi_{j}\right\vert h_{2}\psi_{j}\right)
\left\vert \psi_{j+s}\right\rangle \left\langle \psi_{j+s}\right\vert -\left(
\left.  \psi_{j+s}\right\vert h_{2}\psi_{j}\right)  \left\vert \psi
_{j+s}\right\rangle \left\langle \psi_{j}\right\vert -\left(  \left.  \psi
_{j}\right\vert h_{2}\psi_{j+s}\right)  \left\vert \psi_{j}\right\rangle
\left\langle \psi_{j+s}\right\vert \right\} \nonumber\\
&  \qquad+\sum_{1\leq j,k\leq s}n_{j}^{\frac{1}{2}}n_{k}^{\frac{1}{2}%
}e^{i\left(  \theta_{j}-\theta_{k}\right)  }S_{N-1}\left(  \jmath k\right)
\left\{  \left(  \left.  \psi_{j+s}\right\vert h_{2}\psi_{k+s}\right)
\left\vert \psi_{k}\right\rangle \left\langle \psi_{j}\right\vert +\left(
\left.  \psi_{j}\right\vert h_{2}\psi_{k}\right)  \left\vert \psi
_{k+s}\right\rangle \left\langle \psi_{j+s}\right\vert \right. \nonumber\\
&  \qquad\qquad\qquad\qquad\qquad\qquad\left.  -\left(  \left.  \psi
_{j+s}\right\vert h_{2}\psi_{k}\right)  \left\vert \psi_{k+s}\right\rangle
\left\langle \psi_{j}\right\vert -\left(  \left.  \psi_{j}\right\vert
h_{2}\psi_{k+s}\right)  \left\vert \psi_{k}\right\rangle \left\langle
\psi_{j+s}\right\vert \right\} \nonumber\\
&  \qquad+2\sum_{1\leq j,k\leq2s}n_{j}n_{k}S_{N-2}\left(  \jmath k\right)
\left\{  \left(  \left.  \psi_{k}\right\vert h_{2}\psi_{k}\right)  \left\vert
\psi_{j}\right\rangle \left\langle \psi_{j}\right\vert -\left(  \left.
\psi_{k}\right\vert h_{2}\psi_{j}\right)  \left\vert \psi_{k}\right\rangle
\left\langle \psi_{j}\right\vert \right\}  \left.  \overset{}{}\right\}
\left\vert \psi_{\mu}\right\rangle . \label{der1}%
\end{align}
The derivatives wrt $\mathbf{\psi}$ are
\begin{align}
&  \frac{\delta\mathcal{H}\left(  \mathbf{\psi,n}\right)  }{\delta\psi_{\mu}%
}=\left\langle \psi_{\mu}\right\vert \left\{  \sum_{1\leq j\leq s}\left\{
N_{j}\left\{  \left\vert \psi_{j}\right\rangle \left\langle \psi
_{j}\right\vert +\left\vert \psi_{j+s}\right\rangle \left\langle \psi
_{j+s}\right\vert \right\}  h_{1}+\left(  \left.  \psi_{j+s}\right\vert
h_{2}\psi_{j+s}\right)  \left\vert \psi_{j}\right\rangle \left\langle \psi
_{j}\right\vert \right.  \right. \nonumber\\
&  \qquad\left.  +\left(  \left.  \psi_{j}\right\vert h_{2}\psi_{j}\right)
\left\vert \psi_{j+s}\right\rangle \left\langle \psi_{j+s}\right\vert -\left(
\left.  \psi_{j+s}\right\vert h_{2}\psi_{j}\right)  \left\vert \psi
_{j+s}\right\rangle \left\langle \psi_{j}\right\vert -\left(  \left.  \psi
_{j}\right\vert h_{2}\psi_{j+s}\right)  \left\vert \psi_{j}\right\rangle
\left\langle \psi_{j+s}\right\vert \right\} \nonumber\\
&  \qquad+\sum_{1\leq j,k\leq s}n_{j}^{\frac{1}{2}}n_{k}^{\frac{1}{2}%
}e^{i\left(  \theta_{j}-\theta_{k}\right)  }S_{N-1}\left(  \jmath k\right)
\left\{  \left(  \left.  \psi_{j+s}\right\vert h_{2}\psi_{k+s}\right)
\left\vert \psi_{k}\right\rangle \left\langle \psi_{j}\right\vert +\left(
\left.  \psi_{j}\right\vert h_{2}\psi_{k}\right)  \left\vert \psi
_{k+s}\right\rangle \left\langle \psi_{j+s}\right\vert \right. \nonumber\\
&  \qquad\qquad\qquad\qquad\qquad\qquad\left.  -\left(  \left.  \psi
_{j+s}\right\vert h_{2}\psi_{k}\right)  \left\vert \psi_{k+s}\right\rangle
\left\langle \psi_{j}\right\vert -\left(  \left.  \psi_{j}\right\vert
h_{2}\psi_{k+s}\right)  \left\vert \psi_{k}\right\rangle \left\langle
\psi_{j+s}\right\vert \right\} \nonumber\\
&  \qquad+2\sum_{1\leq j,k\leq2s}n_{j}n_{k}S_{N-2}\left(  \jmath k\right)
\left\{  \left(  \left.  \psi_{k}\right\vert h_{2}\psi_{k}\right)  \left\vert
\psi_{j}\right\rangle \left\langle \psi_{j}\right\vert -\left(  \left.
\psi_{k}\right\vert h_{2}\psi_{j}\right)  \left\vert \psi_{k}\right\rangle
\left\langle \psi_{j}\right\vert \right\}  \left.  \overset{}{}\right\}  ,
\label{der2}%
\end{align}
showing that
\begin{equation}
\frac{\delta\mathcal{H}\left(  \mathbf{\psi,n}\right)  }{\delta\psi_{\mu}%
}=\left(  \frac{\delta\mathcal{H}\left(  \mathbf{\psi,n}\right)  }{\delta
\psi_{\mu}^{\ast}}\right)  ^{\ast}.
\end{equation}
If we define a generalized Fock operator
\begin{align}
&  F\left(  \mathbf{n},\mathbf{\psi},\mathbf{\theta}\right)  =\sum
_{\substack{1\leq j\leq s\\1\leq m\leq2s}}N_{j}\left\{  \left\vert \psi
_{m}\right\rangle \left\langle \psi_{j}\right\vert h_{1mj}+\left\vert
\psi_{j+s}\right\rangle \left\langle \psi_{j+s}\right\vert h_{1mj+s}%
+\left\vert \psi_{m}\right\rangle \left\langle \psi_{j}\right\vert \left(
\left.  \psi_{j+s}\right\vert h_{2}\psi_{j+s}\right)  _{mj}\right. \nonumber\\
&  \qquad\left.  +\,\left\vert \psi_{m}\right\rangle \left\langle \psi
_{j+s}\right\vert \left(  \left.  \psi_{j}\right\vert h_{2}\psi_{j}\right)
_{mj+s}-\left\vert \psi_{m}\right\rangle \left\langle \psi_{j}\right\vert
\left(  \left.  \psi_{j+s}\right\vert h_{2}\psi_{j}\right)  _{mj+s}-\left\vert
\psi_{m}\right\rangle \left\langle \psi_{j+s}\right\vert \left(  \left.
\psi_{j}\right\vert h_{2}\psi_{j+s}\right)  _{mj}\right\} \nonumber\\
&  \qquad+\sum_{\substack{1\leq j,k\leq s\\1\leq m\leq2s}}n_{j}^{\frac{1}{2}%
}n_{k}^{\frac{1}{2}}e^{i\left(  \theta_{j}-\theta_{k}\right)  }S_{N-1}\left(
\jmath k\right)  \left\{  \left\vert \psi_{m}\right\rangle \left\langle
\psi_{j}\right\vert \left(  \left.  \psi_{j+s}\right\vert h_{2}\psi
_{k+s}\right)  _{mk}+\left\vert \psi_{m}\right\rangle \left\langle \psi
_{j+s}\right\vert \left(  \left.  \psi_{j}\right\vert h_{2}\psi_{k}\right)
_{mk+s}\right. \nonumber\\
&  \qquad\qquad\qquad\qquad\qquad\qquad\left.  -\left\vert \psi_{m}%
\right\rangle \left\langle \psi_{j}\right\vert \left(  \left.  \psi
_{j+s}\right\vert h_{2}\psi_{k}\right)  _{mk+s}-\left\vert \psi_{m}%
\right\rangle \left\langle \psi_{j+s}\right\vert \left(  \left.  \psi
_{j}\right\vert h_{2}\psi_{k+s}\right)  _{mk}\right\} \nonumber\\
&  \qquad+2\sum_{\substack{1\leq j,k\leq2s\\1\leq m\leq2s}}n_{j}n_{k}%
S_{N-2}\left(  \jmath k\right)  \left\{  \left\vert \psi_{m}\right\rangle
\left\langle \psi_{j}\right\vert \left(  \left.  \psi_{k}\right\vert h_{2}%
\psi_{k}\right)  _{mj}-\left\vert \psi_{m}\right\rangle \left\langle \psi
_{j}\right\vert \left(  \left.  \psi_{k}\right\vert h_{2}\psi_{j}\right)
_{mk}\right\}  ,
\end{align}
giving%
\begin{equation}
F\left(  \mathbf{n},\mathbf{\psi},\mathbf{\theta}\right)  =\sum_{1\leq
j,m\leq2s}\left\{  N_{j}\left\{  h_{1mj}+A_{mj}-B_{mj}\right\}  +\left\{
C_{mj}-D_{mj}\right\}  +2\overset{}{\left\{  E_{mj}-F_{mj}\right\}  \left\vert
\psi_{m}\right\rangle \left\langle \psi_{j}\right\vert }\right\}  .
\end{equation}
The matrix elements of $\left\{  \mathbb{A},\mathbb{B},\mathbb{C}%
,\mathbb{D},\mathbb{E},\mathbb{F}\right\}  $ for $1\leq m\leq2s$ are given by
\begin{equation}
A_{mj}=\left\{
\begin{array}
[c]{c}%
\left(  \left.  \psi_{j+s}\right\vert h_{2}\psi_{j+s}\right)  _{mj};1\leq
j\leq s\\
\left(  \left.  \psi_{j-s}\right\vert h_{2}\psi_{j-s}\right)  _{mj};s+1\leq
j\leq s
\end{array}
\right.  ,
\end{equation}%
\begin{equation}
B_{mj}=\left\{
\begin{array}
[c]{c}%
\left(  \left.  \psi_{j+s}\right\vert h_{2}\psi_{j}\right)  _{mj+s};1\leq
j\leq s\\
\left(  \left.  \psi_{j-s}\right\vert h_{2}\psi_{j}\right)  _{mj-s};s+1\leq
j\leq s
\end{array}
\right.  ,
\end{equation}
noting that $n_{j+s}=n_{j}$ and $\theta_{j+s}=\theta_{j}$ for $1\leq j\leq s$
we have the phase dependent $2$-electron terms
\begin{equation}
C_{mj}=\left\{
\begin{array}
[c]{c}%
%TCIMACRO{\tsum \limits_{1\leq k\leq s}}%
%BeginExpansion
{\textstyle\sum\limits_{1\leq k\leq s}}
%EndExpansion
n_{j}^{\frac{1}{2}}n_{k}^{\frac{1}{2}}S_{N-1}\left(  \jmath k\right)
e^{i\left(  \theta_{j}-\theta_{k}\right)  }\left(  \left.  \psi_{j+s}%
\right\vert h_{2}\psi_{k+s}\right)  _{mk};1\leq j\leq s\\%
%TCIMACRO{\tsum \limits_{1\leq k\leq s}}%
%BeginExpansion
{\textstyle\sum\limits_{1\leq k\leq s}}
%EndExpansion
n_{j}^{\frac{1}{2}}n_{k}^{\frac{1}{2}}S_{N-1}\left(  \jmath k\right)
e^{i\left(  \theta_{j}-\theta_{k}\right)  }\left(  \left.  \psi_{j-s}%
\right\vert h_{2}\psi_{k}\right)  _{mk+s};s+1\leq j\leq2s
\end{array}
\right.  ,
\end{equation}%
\begin{equation}
D_{mj}=\left\{
\begin{array}
[c]{c}%
%TCIMACRO{\tsum \limits_{1\leq k\leq s}}%
%BeginExpansion
{\textstyle\sum\limits_{1\leq k\leq s}}
%EndExpansion
n_{j}^{\frac{1}{2}}n_{k}^{\frac{1}{2}}S_{N-1}\left(  \jmath k\right)
e^{i\left(  \theta_{j}-\theta_{k}\right)  }\left(  \left.  \psi_{j+s}%
\right\vert h_{2}\psi_{k}\right)  _{mk+s};1\leq j\leq s\\%
%TCIMACRO{\tsum \limits_{1\leq k\leq s}}%
%BeginExpansion
{\textstyle\sum\limits_{1\leq k\leq s}}
%EndExpansion
n_{j}^{\frac{1}{2}}n_{k}^{\frac{1}{2}}S_{N-1}\left(  \jmath k\right)
e^{i\left(  \theta_{j}-\theta_{k}\right)  }\left(  \left.  \psi_{j-s}%
\right\vert h_{2}\psi_{k}\right)  _{mk+s};s+1\leq j\leq2s
\end{array}
\right.  ,
\end{equation}
the phase independent $2$-electron terms
\begin{equation}
F_{mj}=\sum_{1\leq k\leq2s}n_{j}n_{k}S_{N-2}\left(  \jmath k\right)  \left(
\left.  \psi_{k}\right\vert h_{2}\psi_{j}\right)  _{mk};1\leq j\leq2s,
\end{equation}
and%
\begin{equation}
E_{mj}=\sum_{1\leq k\leq2s}n_{j}n_{k}S_{N-2}\left(  \jmath k\right)  \left(
\left.  \psi_{k}\right\vert h_{2}\psi_{k}\right)  _{mj};1\leq j\leq2s.
\end{equation}
The derivatives eqs. $\left(  \ref{der1}\right)  $ and eqs. $\left(
\ref{der2}\right)  $ can be expressed as
\begin{align}
\frac{\delta\mathcal{H}\left(  \mathbf{\psi,n}\right)  }{\delta\psi_{\mu
}^{\ast}}  &  =F\left(  \mathbf{n},\mathbf{\psi},\mathbf{\theta}\right)
\left\vert \psi_{\mu}\right\rangle \\
\frac{\delta\mathcal{H}\left(  \mathbf{\psi,n}\right)  }{\delta\psi_{\mu}}  &
=\left\langle \psi_{\mu}\right\vert F\left(  \mathbf{n},\mathbf{\psi
},\mathbf{\theta}\right)  .
\end{align}


\section{Fock Operator and Energy}%

\begin{equation}
f\left(  \mathbf{x}\right)  =f_{0}+f_{1}\left(  \mathbf{x}\right)
+\cdots+f_{n}\left(  \mathbf{x}\right)  =%
%TCIMACRO{\tsum \limits_{0\leq m\leq n}}%
%BeginExpansion
{\textstyle\sum\limits_{0\leq m\leq n}}
%EndExpansion
f_{m}\left(  \mathbf{x}\right)
\end{equation}
where
\begin{equation}
f_{m}\left(  \mathbf{x}\right)  =%
%TCIMACRO{\tsum \limits_{1\leq j_{1},\cdots,j_{m}\leq r}}%
%BeginExpansion
{\textstyle\sum\limits_{1\leq j_{1},\cdots,j_{m}\leq r}}
%EndExpansion
c_{j_{1}\cdots j_{m}}x_{j_{1}}\cdots x_{j_{m}}.
\end{equation}
Consider
\begin{align}
\frac{\partial f_{m}}{\partial x_{l}}\left(  \mathbf{x}\right)   &  =%
%TCIMACRO{\tsum \limits_{1\leq j_{1},\cdots,j_{m-1}\neq l\leq r}}%
%BeginExpansion
{\textstyle\sum\limits_{1\leq j_{1},\cdots,j_{m-1}\neq l\leq r}}
%EndExpansion
c_{j_{1}\cdots j_{l}\cdots j_{m}}x_{j_{1}}\cdots x_{j_{m-1}}\\
x_{l}\frac{\partial f_{m}}{\partial x_{l}}\left(  \mathbf{x}\right)   &  =%
%TCIMACRO{\tsum \limits_{1\leq j_{1},\cdots,j_{m-1}\neq l\leq r}}%
%BeginExpansion
{\textstyle\sum\limits_{1\leq j_{1},\cdots,j_{m-1}\neq l\leq r}}
%EndExpansion
c_{j_{1}\cdots j_{l}\cdots j_{m}}x_{j_{1}}\cdots x_{j_{m-1}}x_{l}\\%
%TCIMACRO{\tsum \limits_{1\leq l\leq r}}%
%BeginExpansion
{\textstyle\sum\limits_{1\leq l\leq r}}
%EndExpansion
x_{l}\frac{\partial f_{m}}{\partial x_{l}}\left(  \mathbf{x}\right)   &
=mf_{m}\left(  \mathbf{x}\right)
\end{align}
hence%
\begin{align}%
%TCIMACRO{\tsum \limits_{1\leq l\leq r}}%
%BeginExpansion
{\textstyle\sum\limits_{1\leq l\leq r}}
%EndExpansion
x_{l}\frac{\partial f}{\partial x_{l}}\left(  \mathbf{x}\right)   &  =%
%TCIMACRO{\tsum \limits_{1\leq l\leq r}}%
%BeginExpansion
{\textstyle\sum\limits_{1\leq l\leq r}}
%EndExpansion
x_{l}\frac{\partial}{\partial x_{l}}%
%TCIMACRO{\tsum \limits_{0\leq m\leq n}}%
%BeginExpansion
{\textstyle\sum\limits_{0\leq m\leq n}}
%EndExpansion
f_{m}\left(  \mathbf{x}\right)  =%
%TCIMACRO{\tsum \limits_{1\leq l\leq r}}%
%BeginExpansion
{\textstyle\sum\limits_{1\leq l\leq r}}
%EndExpansion%
%TCIMACRO{\tsum \limits_{0\leq m\leq n}}%
%BeginExpansion
{\textstyle\sum\limits_{0\leq m\leq n}}
%EndExpansion
x_{l}\frac{\partial f_{m}\left(  \mathbf{x}\right)  }{\partial x_{l}}=%
%TCIMACRO{\tsum \limits_{0\leq m\leq n}}%
%BeginExpansion
{\textstyle\sum\limits_{0\leq m\leq n}}
%EndExpansion%
%TCIMACRO{\tsum \limits_{1\leq l\leq r}}%
%BeginExpansion
{\textstyle\sum\limits_{1\leq l\leq r}}
%EndExpansion
x_{l}\frac{\partial f_{m}\left(  \mathbf{x}\right)  }{\partial x_{l}}=%
%TCIMACRO{\tsum \limits_{0\leq m\leq n}}%
%BeginExpansion
{\textstyle\sum\limits_{0\leq m\leq n}}
%EndExpansion
mf_{m}\left(  \mathbf{x}\right) \\
&  =f_{1}\left(  \mathbf{x}\right)  +2f_{2}\left(  \mathbf{x}\right)
+\cdots+nf_{n}\left(  \mathbf{x}\right)  .
\end{align}
Thus
\begin{equation}%
%TCIMACRO{\tsum \limits_{1\leq l\leq2s}}%
%BeginExpansion
{\textstyle\sum\limits_{1\leq l\leq2s}}
%EndExpansion
\left\langle \left.  \psi_{l}\right\vert F\left(  \mathbf{\psi,n}\right)
\psi_{l}\right\rangle =Tr\left\{  F\left(  \mathbf{\psi,n}\right)  \right\}
=\mathcal{H}_{1}\left(  \mathbf{\psi,n}\right)  +2\mathcal{H}_{2}\left(
\mathbf{\psi,n}\right)
\end{equation}
showing that
\begin{equation}
Tr\left\{  F\left(  \mathbf{\psi,n}\right)  \right\}  -\mathcal{H}\left(
\mathbf{\psi,n}\right)  =\mathcal{H}_{2}\left(  \mathbf{\psi,n}\right)
\end{equation}
and
\begin{equation}
2\mathcal{H}\left(  \mathbf{\psi,n}\right)  =Tr\left\{  F\left(
\mathbf{\psi,n}\right)  \right\}  +\mathcal{H}_{1}\left(  \mathbf{\psi
,n}\right)
\end{equation}
which gives
\begin{equation}
\mathcal{H}\left(  \mathbf{\psi,n}\right)  =Tr\left\{  F\left(  \mathbf{\psi
,n}\right)  \right\}  -\mathcal{H}_{2}\left(  \mathbf{\psi,n}\right)
\end{equation}
and
\begin{equation}
\mathcal{H}\left(  \mathbf{\psi,n}\right)  =\frac{1}{2}\left\{  Tr\left\{
F\left(  \mathbf{\psi,n}\right)  \right\}  +\mathcal{H}_{1}\left(
\mathbf{\psi,n}\right)  \right\}  .
\end{equation}
Noting that
\begin{equation}
Tr\left\{  F\left(  \mathbf{\psi,n}\right)  \right\}  =%
%TCIMACRO{\tsum \limits_{1\leq l\leq2s}}%
%BeginExpansion
{\textstyle\sum\limits_{1\leq l\leq2s}}
%EndExpansion
\varepsilon_{l}\left(  \mathbf{\psi,n}\right)
\end{equation}
and
\begin{equation}
\mathcal{H}_{1}\left(  \mathbf{\psi,n}\right)  =%
%TCIMACRO{\tsum \limits_{1\leq l\leq2s}}%
%BeginExpansion
{\textstyle\sum\limits_{1\leq l\leq2s}}
%EndExpansion
N_{l}\left(  \left.  \psi_{l}\right\vert h_{1}\psi_{l}\right)
\end{equation}
and
\begin{equation}
\left(  \left.  \psi_{l}\right\vert h_{1}\psi_{l}\right)  \leq0,
\end{equation}
it seems that the CGSO's should be occupied by an aufbau principle based on
$\left\{  \left(  \left.  \psi_{l}\right\vert h_{1}\psi_{l}\right)  \right\}
,$ i.e.
\begin{equation}
N_{1}\geq N_{2}\geq\cdots\geq N_{s}%
\end{equation}
corresponding to
\begin{equation}
\left(  \left.  \psi_{1}\right\vert h_{1}\psi_{1}\right)  \leq\left(  \left.
\psi_{1+s}\right\vert h_{1}\psi_{1+s}\right)  \leq\left(  \left.  \psi
_{2}\right\vert h_{1}\psi_{2}\right)  \leq\left(  \left.  \psi_{2+s}%
\right\vert h_{1}\psi_{2+s}\right)  \leq\cdots\leq\left(  \left.  \psi
_{s}\right\vert h_{1}\psi_{s}\right)  \leq\left(  \left.  \psi_{2s}\right\vert
h_{1}\psi_{2s}\right)
\end{equation}
Consider the


\end{document}